% Wiederverwendbare Nullproduktprofile; Einbindung am Ende von Band 33.
\chapter{Nullproduktprofile}

Die Nullprodukte einer Halbgruppe lassen sich bereits in ihrer
Potenzhalbgruppe erkennen. Dabei werden keine Einermengen erkannt:
zunächst erhält man Klassen mit gleichem Verhalten gegenüber allen
linken und rechten Nulltests. Bei endlichen Trägern bestimmen die Größen
unterer Mengen, wie viele Grundelemente in jeder dieser Klassen liegen.

Wir schreiben \(K_S=\PowNE S\), \(z_S=\{0_S\}\) und entsprechend
\(K_T=\PowNE T\), \(z_T=\{0_T\}\). Nebeneinanderschreiben bezeichnet
bei Grundelementen die jeweilige Halbgruppenoperation und bei Teilmengen
das Mengenprodukt aus \FormulaRefAuto{PowerSemigroupProductDef}[def].
Die Abkürzungen
\[
 \begin{gathered}
 H_S=\SemigroupWithZeroAlg{S}{\star}{0_S},\qquad
 H_T=\SemigroupWithZeroAlg{T}{\diamond}{0_T},\\
 I=\SemigroupIso{\Phi}{K_S,\star_{\mathcal P}}{K_T,\diamond_{\mathcal P}}.
 \end{gathered}
\]
werden in Aussagen und Tabellen nur zur Kürzung wiederholter
Voraussetzungen verwendet. \(I\) enthält insbesondere die Bijektivität
von \(\Phi\).

\Needspace{16\baselineskip}
\section{Einermengen und Nullprodukte}

\Needspace{9\baselineskip}
\FormulaThmDeltaK[Einermengen sind nichtleere Teilmengen]{%
 a\in S\vdash\{a\}\in\PowNE S
}{SingletonPowerMember}{\DeltaRow{Mengen}{S\dsep a}}
\begin{tabproof}
 \proofstep{1}{a\in S}{\rA}
 \proofstep{1}{\{a\}\subseteq S}{\FormulaRefAuto{a\in A\vdash\{a\}\subseteq A}[thm]{1}}
 \proofstep{}{a\in\{a\}}{\FormulaRefAuto{a\in\{a\}}[thm]}
 \proofstep{}{\{a\}\neq\varnothing}{\FormulaRefAuto{a\in A\vdash A\neq\varnothing}[thm]{3}}
 \proofstep{1}{\{a\}\in\PowNE S}{\FormulaRefAuto{NonemptyPowersetMembership}[thm]{\rAI{2,4}}}
\end{tabproof}

\Needspace{9\baselineskip}
\FormulaThmDeltaK[Konstante Produkte nichtleerer Teilmengen]{%
 \SemigroupAlg{S}{\star}\dsep A,B\in K_S\vdash
 (AB=\{d\}\leftrightarrow\forall a\in A\,\forall b\in B\;(ab=d))
}{PowerSemigroupConstantProduct}{%
 \DeltaRow{Mengen}{S\dsep A\dsep B\dsep a\dsep b\dsep d\dsep w}}
\Needspace{6\baselineskip}
\begin{tabproofsplitwide}
\Needspace{5\baselineskip}
\proofpartwide{Hinrichtung}
 \proofstepwidestar[1]{\SemigroupAlg{S}{\star}}{\rA}
 \proofstepwidestar[2]{A,B\in K_S}{\rA}
 \proofstepwidestar[3]{AB=\{d\}}{\rA}
 \proofstepwidestar[4]{a\in A}{\rA}
 \proofstepwidestar[5]{b\in B}{\rA}
 \proofstepwidestar[1,2,4,5]{ab\in AB}{\FormulaRefAuto{PowerSemigroupProductContains}[thm]{1,2,4,5}}
 \proofstepwidestar[1,2,3,4,5]{ab\in\{d\}}{\rIE{3,6}}
 \proofstepwidestar[1,2,3,4,5]{ab=d}{\FormulaRefAuto{x\in\{a\}\eqvdash x=a}[thm]{7}}
 \proofstepwidestar[1,2,3]{\forall a\in A\,\forall b\in B\;(ab=d)}{\rUI{\rRI{4,\rUI{\rRI{5,8}}}}}
 \proofstepwidestar[1,2]{AB=\{d\}\to\forall a\in A\,\forall b\in B\;(ab=d)}{\rRI{3,9}}
\closeproofpartwide
\Needspace{5\baselineskip}
\proofpartwide{Rückrichtung}
 \setcounter{proofstepnr}{10}
 \proofstepwidestar[11]{\forall a\in A\,\forall b\in B\;(ab=d)}{\rA}
 \proofstepwidestar[12]{w\in AB}{\rA}
 \proofstepwidestar[1,2,12]{\exists a\in A\,\exists b\in B\;(w=ab)}{\FormulaRefAuto{PowerSemigroupProductWitnesses}[thm]{1,2,12}}
 \proofstepwidestar[14]{a\in A\land b\in B\land w=ab}{\rA}
 \proofstepwidestar[11,14]{ab=d}{\rRE{\rUE{\rRE{\rUE{11},14}},14}}
 \proofstepwidestar[11,14]{w=d}{\rChain{\rAEb{\rAEb{14}},15}}
 \proofstepwidestar[11,14]{w\in\{d\}}{\FormulaRefAuto{x\in\{a\}\eqvdash x=a}[thm]{16}}
 \proofstepwidestar[1,2,11,12]{w\in\{d\}}{\rEEStar{13,14\text{--}17}}
 \proofstepwidestar[1,2,11]{AB\subseteq\{d\}}{\FormulaRefAuto{SubsetIntroductionRule}[thm]{[12]\vdots18}}
 \proofstepwidestar[1,2]{AB\neq\varnothing}{\FormulaRefAuto{PowerSemigroupProductNonempty}[thm]{1,2}}
 \proofstepwidestar[1,2]{\exists w\;w\in AB}{\FormulaRefAuto{A\neq\varnothing\eqvdash\exists x\,(x\in A)}[thm]{20}}
 \proofstepwidestar[22]{w\in AB}{\rA}
 \proofstepwidestar[1,2,11,22]{w\in\{d\}}{\FormulaRefAuto{A\subseteq B\dsep x\in A\vdash x\in B}[thm]{19,22}}
 \proofstepwidestar[1,2,11,22]{w=d}{\FormulaRefAuto{x\in\{a\}\eqvdash x=a}[thm]{23}}
 \proofstepwidestar[1,2,11,22]{d\in AB}{\rIE{24,22}}
 \proofstepwidestar[1,2,11]{d\in AB}{\rEE{21,[22]\vdots25}}
 \proofstepwidestar[1,2,11]{\{d\}\subseteq AB}{\FormulaRefAuto{a\in A\vdash\{a\}\subseteq A}[thm]{26}}
 \proofstepwidestar[1,2,11]{AB=\{d\}}{\FormulaRefAuto{A\subseteq B,B\subseteq A\vdash A=B}[thm]{19,27}}
 \proofstepwidestar[1,2]{(\forall a\in A\,\forall b\in B\;(ab=d))\to AB=\{d\}}{\rRI{11,28}}
 \proofstepwidestar[1,2]{AB=\{d\}\leftrightarrow\forall a\in A\,\forall b\in B\;(ab=d)}{\rLRI{10,29}}
\closeproofpartwide
\end{tabproofsplitwide}

\Needspace{9\baselineskip}
\FormulaThmDeltaK[Produkt zweier Einermengen]{%
 \SemigroupAlg{S}{\star}\dsep a,b\in S\vdash\{a\}\{b\}=\{ab\}
}{SingletonPowerProduct}{\DeltaRow{Mengen}{S\dsep a\dsep b\dsep x\dsep y}}
\begin{tabproofwide}
 \proofstepwidestar[1]{\SemigroupAlg{S}{\star}}{\rA}
 \proofstepwidestar[2]{a,b\in S}{\rA}
 \proofstepwidestar[2]{\{a\},\{b\}\in K_S}{\FormulaRefAuto{SingletonPowerMember}[thm]{2}}
 \proofstepwidestar[4]{x\in\{a\}}{\rA}
 \proofstepwidestar[5]{y\in\{b\}}{\rA}
 \proofstepwidestar[4]{x=a}{\FormulaRefAuto{x\in\{a\}\eqvdash x=a}[thm]{4}}
 \proofstepwidestar[5]{y=b}{\FormulaRefAuto{x\in\{a\}\eqvdash x=a}[thm]{5}}
 \proofstepwidestar[4,5]{xy=ab}{\rIE{7,\rIE{6,\rII}}}
 \proofstepwidestar[]{\forall x\in\{a\}\,\forall y\in\{b\}\;(xy=ab)}{\rUI{\rRI{4,\rUI{\rRI{5,8}}}}}
 \proofstepwidestar[1,2]{\{a\}\{b\}=\{ab\}}{\FormulaRefAuto{PowerSemigroupConstantProduct}[thm]{1,3,9}}
\end{tabproofwide}

\Needspace{9\baselineskip}
\FormulaThmDeltaK[Die nichtleere Potenzhalbgruppe hat eine Null]{%
 H_S\vdash\SemigroupWithZeroAlg{K_S}{\star_{\mathcal P}}{z_S}
}{PowerSemigroupWithZero}{\DeltaRow{Mengen}{S\dsep0_S\dsep A\dsep x}}
\begin{tabproofwide}
 \proofstepwidestar[1]{H_S}{\rA}
 \proofstepwidestar[1]{\SemigroupAlg{S}{\star}}{\FormulaRefAuto{SemigroupWithZeroBaseAxiom}[ax]{1}}
 \proofstepwidestar[1]{0_S\in S}{\FormulaRefAuto{SemigroupWithZeroCarrierAxiom}[ax]{1}}
 \proofstepwidestar[1]{\SemigroupAlg{K_S}{\star_{\mathcal P}}}{\FormulaRefAuto{NonemptyPowerSemigroup}[thm]{2}}
 \proofstepwidestar[1]{z_S\in K_S}{\FormulaRefAuto{SingletonPowerMember}[thm]{3}}
 \proofstepwidestar[6]{A\in K_S}{\rA}
 \proofstepwidestar[7]{x\in A}{\rA}
 \proofstepwidestar[6,7]{x\in S}{\FormulaRefAuto{PowerSemigroupCarrierElement}[thm]{6,7}}
 \proofstepwidestar[1,6,7]{0_Sx=0_S\land x0_S=0_S}{\rAI{\FormulaRefAuto{SemigroupWithZeroLeftAbsorptionAxiom}[ax]{1,8},\FormulaRefAuto{SemigroupWithZeroRightAbsorptionAxiom}[ax]{1,8}}}
 \proofstepwidestar[10]{u\in z_S}{\rA}
 \proofstepwidestar[10]{u=0_S}{\FormulaRefAuto{x\in\{a\}\eqvdash x=a}[thm]{10}}
 \proofstepwidestar[1,6,7,10]{ux=0_S\land xu=0_S}{\rIE{\FormulaRefAuto{a=b\vdash b=a}[thm]{11},9}}
 \proofstepwidestar[1,6]{\forall u\in z_S\,\forall x\in A\;(ux=0_S)}{\rUI{\rRI{10,\rUI{\rRI{7,\rAEa{12}}}}}}
 \proofstepwidestar[1,6]{\forall x\in A\,\forall u\in z_S\;(xu=0_S)}{\rUI{\rRI{7,\rUI{\rRI{10,\rAEb{12}}}}}}
 \proofstepwidestar[1,6]{z_SA=z_S}{\FormulaRefAuto{PowerSemigroupConstantProduct}[thm]{2,5,6,13}}
 \proofstepwidestar[1,6]{Az_S=z_S}{\FormulaRefAuto{PowerSemigroupConstantProduct}[thm]{2,6,5,14}}
 \proofstepwidestar[1]{\SemigroupWithZeroAlg{K_S}{\star_{\mathcal P}}{z_S}}{\FormulaRefAuto{SemigroupWithZeroDef}[def]{4,5,\rUI{\rRI{6,15}},\rUI{\rRI{6,16}}}}
\end{tabproofwide}

\Needspace{9\baselineskip}
\FormulaThmDeltaK[Isomorphismen erhalten die Null der Potenzhalbgruppe]{%
 H_S\dsep H_T\dsep I\vdash\Phi(z_S)=z_T
}{PowerSemigroupZeroIsoValue}{\DeltaRow{Mengen}{S\dsep T\dsep0_S\dsep0_T\dsep Y}}
\begin{tabproofwide}
 \proofstepwidestar[1]{H_S}{\rA}
 \proofstepwidestar[2]{H_T}{\rA}
 \proofstepwidestar[3]{I}{\rA}
 \proofstepwidestar[1]{\SemigroupWithZeroAlg{K_S}{\star_{\mathcal P}}{z_S}}{\FormulaRefAuto{PowerSemigroupWithZero}[thm]{1}}
 \proofstepwidestar[2]{\SemigroupWithZeroAlg{K_T}{\diamond_{\mathcal P}}{z_T}}{\FormulaRefAuto{PowerSemigroupWithZero}[thm]{2}}
 \proofstepwidestar[1]{z_S\in K_S}{\FormulaRefAuto{SemigroupWithZeroCarrierAxiom}[ax]{4}}
 \proofstepwidestar[3]{\Phi:K_S\to K_T}{\FormulaRefAuto{SemigroupIsoFunction}[thm]{3}}
 \proofstepwidestar[1,3]{\Phi(z_S)\in K_T}{\FormulaRefAuto{F\colon A\to B\dsep x\in A\vdash F(x)\in B}[thm]{7,6}}
 \proofstepwidestar[1]{\forall X\in K_S\;(z_SX=z_S\land Xz_S=z_S)}{\FormulaRefAuto{SemigroupWithZeroDef}[def]{4}}
 \proofstepwidestar[1,3]{\forall Y\in K_T\;(\Phi(z_S)Y=\Phi(z_S)\land Y\Phi(z_S)=\Phi(z_S))}{\FormulaRefAuto{SemigroupIsoIdentityZeroConditions}[thm]{3,\rAI{6,6},9}}
 \proofstepwidestar[3]{\SemigroupAlg{K_T}{\diamond_{\mathcal P}}}{\FormulaRefAuto{SemigroupIsoTarget}[thm]{3}}
 \proofstepwidestar[1,3]{\SemigroupWithZeroAlg{K_T}{\diamond_{\mathcal P}}{\Phi(z_S)}}{\FormulaRefAuto{SemigroupWithZeroDef}[def]{11,8,10}}
 \proofstepwidestar[1,2,3]{\Phi(z_S)=z_T}{\FormulaRefAuto{SemigroupZeroUnique}[thm]{12,5}}
\end{tabproofwide}

\Needspace{9\baselineskip}
\FormulaThmDeltaK[Linke Zerlegung eines konstanten Mengenprodukts]{%
 \SemigroupAlg{S}{\star}\dsep A,X\in K_S\vdash
 (AX=\{d\}\leftrightarrow\forall a\in A\;\{a\}X=\{d\})
}{PowerSemigroupLeftAtomization}{\DeltaRow{Mengen}{S\dsep A\dsep X\dsep a\dsep x\dsep u\dsep d}}
\Needspace{6\baselineskip}
\begin{tabproofsplitwide}
\Needspace{5\baselineskip}
\proofpartwide{Hinrichtung}
 \proofstepwidestar[1]{\SemigroupAlg{S}{\star}}{\rA}
 \proofstepwidestar[2]{A,X\in K_S}{\rA}
 \proofstepwidestar[3]{AX=\{d\}}{\rA}
 \proofstepwidestar[1,2]{AX=\{d\}\leftrightarrow \forall a\in A\,\forall x\in X\;(ax=d)}{\FormulaRefAuto{PowerSemigroupConstantProduct}[thm]{1,2}}
 \proofstepwidestar[1,2,3]{\forall a\in A\,\forall x\in X\;(ax=d)}{\FormulaRefAuto{P\leftrightarrow Q,P\vdash Q}{4,3}}
 \proofstepwidestar[6]{a\in A}{\rA}
 \proofstepwidestar[7]{u\in\{a\}}{\rA}
 \proofstepwidestar[8]{x\in X}{\rA}
 \proofstepwidestar[7]{u=a}{\FormulaRefAuto{x\in\{a\}\eqvdash x=a}[thm]{7}}
 \proofstepwidestar[1,2,3,6,8]{ax=d}{\rRE{\rUE{\rRE{\rUE{5},6}},8}}
 \proofstepwidestar[1,2,3,6,7,8]{ux=d}{\rIE{\FormulaRefAuto{a=b\vdash b=a}[thm]{9},10}}
 \proofstepwidestar[1,2,3,6]{\forall u\in\{a\}\,\forall x\in X\;(ux=d)}{\rUI{\rRI{7,\rUI{\rRI{8,11}}}}}
 \proofstepwidestar[2,6]{a\in S}{\FormulaRefAuto{PowerSemigroupCarrierElement}[thm]{2,6}}
 \proofstepwidestar[2,6]{\{a\}\in K_S}{\FormulaRefAuto{SingletonPowerMember}[thm]{13}}
 \proofstepwidestar[1,2,3,6]{\{a\}X=\{d\}}{\FormulaRefAuto{PowerSemigroupConstantProduct}[thm]{1,2,14,12}}
 \proofstepwidestar[1,2,3]{\forall a\in A\;\{a\}X=\{d\}}{\rUI{\rRI{6,15}}}
 \proofstepwidestar[1,2]{AX=\{d\}\to\forall a\in A\;\{a\}X=\{d\}}{\rRI{3,16}}
\closeproofpartwide
\Needspace{5\baselineskip}
\proofpartwide{Rückrichtung}
 \setcounter{proofstepnr}{17}
 \proofstepwidestar[18]{a\in A}{\rA}
 \proofstepwidestar[19]{x\in X}{\rA}
 \proofstepwidestar[20]{\forall a\in A\;\{a\}X=\{d\}}{\rA}
 \proofstepwidestar[18,20]{\{a\}X=\{d\}}{\rRE{\rUE{20},18}}
 \proofstepwidestar[2,18]{\{a\}\in K_S}{\FormulaRefAuto{SingletonPowerMember}[thm]{\FormulaRefAuto{PowerSemigroupCarrierElement}[thm]{2,18}}}
 \proofstepwidestar[]{a\in\{a\}}{\FormulaRefAuto{a\in\{a\}}[thm]}
 \proofstepwidestar[1,2,18,19,20]{ax=d}{\FormulaRefAuto{PowerSemigroupConstantProduct}[thm]{1,2,22,21,23,19}}
 \proofstepwidestar[1,2,20]{\forall a\in A\,\forall x\in X\;(ax=d)}{\rUI{\rRI{18,\rUI{\rRI{19,24}}}}}
 \proofstepwidestar[1,2,20]{AX=\{d\}}{\FormulaRefAuto{PowerSemigroupConstantProduct}[thm]{1,2,25}}
 \proofstepwidestar[1,2]{(\forall a\in A\;\{a\}X=\{d\})\to AX=\{d\}}{\rRI{20,26}}
 \proofstepwidestar[1,2]{AX=\{d\}\leftrightarrow\forall a\in A\;\{a\}X=\{d\}}{\rLRI{17,27}}
\closeproofpartwide
\end{tabproofsplitwide}


\Needspace{9\baselineskip}
\FormulaThmDeltaK[Rechte Zerlegung eines konstanten Mengenprodukts]{%
 \SemigroupAlg{S}{\star}\dsep A,X\in K_S\vdash
 (XA=\{d\}\leftrightarrow\forall a\in A\;X\{a\}=\{d\})
}{PowerSemigroupRightAtomization}{\DeltaRow{Mengen}{S\dsep A\dsep X\dsep a\dsep x\dsep u\dsep d}}
\Needspace{6\baselineskip}
\begin{tabproofsplitwide}
\Needspace{5\baselineskip}
\proofpartwide{Hinrichtung}
 \proofstepwidestar[1]{\SemigroupAlg{S}{\star}}{\rA}
 \proofstepwidestar[2]{A,X\in K_S}{\rA}
 \proofstepwidestar[3]{XA=\{d\}}{\rA}
 \proofstepwidestar[1,2]{XA=\{d\}\leftrightarrow \forall x\in X\,\forall a\in A\;(xa=d)}{\FormulaRefAuto{PowerSemigroupConstantProduct}[thm]{1,2}}
 \proofstepwidestar[1,2,3]{\forall x\in X\,\forall a\in A\;(xa=d)}{\FormulaRefAuto{P\leftrightarrow Q,P\vdash Q}{4,3}}
 \proofstepwidestar[6]{a\in A}{\rA}
 \proofstepwidestar[7]{u\in\{a\}}{\rA}
 \proofstepwidestar[8]{x\in X}{\rA}
 \proofstepwidestar[7]{u=a}{\FormulaRefAuto{x\in\{a\}\eqvdash x=a}[thm]{7}}
 \proofstepwidestar[1,2,3,6,8]{xa=d}{\rRE{\rUE{\rRE{\rUE{5},8}},6}}
 \proofstepwidestar[1,2,3,6,7,8]{xu=d}{\rIE{\FormulaRefAuto{a=b\vdash b=a}[thm]{9},10}}
 \proofstepwidestar[1,2,3,6]{\forall x\in X\,\forall u\in\{a\}\;(xu=d)}{\rUI{\rRI{8,\rUI{\rRI{7,11}}}}}
 \proofstepwidestar[2,6]{a\in S}{\FormulaRefAuto{PowerSemigroupCarrierElement}[thm]{2,6}}
 \proofstepwidestar[2,6]{\{a\}\in K_S}{\FormulaRefAuto{SingletonPowerMember}[thm]{13}}
 \proofstepwidestar[1,2,3,6]{X\{a\}=\{d\}}{\FormulaRefAuto{PowerSemigroupConstantProduct}[thm]{1,2,14,12}}
 \proofstepwidestar[1,2,3]{\forall a\in A\;X\{a\}=\{d\}}{\rUI{\rRI{6,15}}}
 \proofstepwidestar[1,2]{XA=\{d\}\to\forall a\in A\;X\{a\}=\{d\}}{\rRI{3,16}}
\closeproofpartwide
\Needspace{5\baselineskip}
\proofpartwide{Rückrichtung}
 \setcounter{proofstepnr}{17}
 \proofstepwidestar[18]{a\in A}{\rA}
 \proofstepwidestar[19]{x\in X}{\rA}
 \proofstepwidestar[20]{\forall a\in A\;X\{a\}=\{d\}}{\rA}
 \proofstepwidestar[18,20]{X\{a\}=\{d\}}{\rRE{\rUE{20},18}}
 \proofstepwidestar[2,18]{\{a\}\in K_S}{\FormulaRefAuto{SingletonPowerMember}[thm]{\FormulaRefAuto{PowerSemigroupCarrierElement}[thm]{2,18}}}
 \proofstepwidestar[]{a\in\{a\}}{\FormulaRefAuto{a\in\{a\}}[thm]}
 \proofstepwidestar[1,2,18,19,20]{xa=d}{\FormulaRefAuto{PowerSemigroupConstantProduct}[thm]{1,2,22,21,23,19}}
 \proofstepwidestar[1,2,20]{\forall x\in X\,\forall a\in A\;(xa=d)}{\rUI{\rRI{19,\rUI{\rRI{18,24}}}}}
 \proofstepwidestar[1,2,20]{XA=\{d\}}{\FormulaRefAuto{PowerSemigroupConstantProduct}[thm]{1,2,25}}
 \proofstepwidestar[1,2]{(\forall a\in A\;X\{a\}=\{d\})\to XA=\{d\}}{\rRI{20,26}}
 \proofstepwidestar[1,2]{XA=\{d\}\leftrightarrow\forall a\in A\;X\{a\}=\{d\}}{\rLRI{17,27}}
\closeproofpartwide
\end{tabproofsplitwide}

\Needspace{16\baselineskip}
\section{Präordnung und Quotientenordnung}

\Needspace{12\baselineskip}
\FormulaDefDeltaK[Vergleich der Nullproduktprofile]{%
 A\preceq_S B\coloneqq
 \forall X\in K_S\;((BX=z_S\to AX=z_S)\land(XB=z_S\to XA=z_S))
}{NullProfilePreorderDef}{%
 \DeltaPrem{Halbgruppe mit Null}{H_S}
 \DeltaRow{Mengen}{A\dsep B\dsep X\dsep S\dsep0_S}}

\Needspace{12\baselineskip}
\FormulaDefDeltaK[Gleiches Nullproduktprofil]{%
 A\sim_S B\coloneqq(A\preceq_S B\land B\preceq_S A)
}{NullProfileEquivalenceDef}{%
 \DeltaPrem{Halbgruppe mit Null}{H_S}
 \DeltaRow{Mengen}{A\dsep B\dsep S\dsep0_S}}

\Needspace{9\baselineskip}
\FormulaThmDeltaK[Nullproduktprofile bilden eine Präordnung]{%
 H_S\vdash\operatorname{PrO}(K_S,\preceq_S)
}{NullProfilePreorder}{\DeltaRow{Mengen}{S\dsep A\dsep B\dsep C\dsep X}}
\Needspace{6\baselineskip}
\begin{tabproofsplitwide}
\Needspace{8\baselineskip}
\proofpartwideindR[Reflexivität]{H_S\dsep A\in K_S\vdash A\preceq_S A}{H_S\dsep A\in K_S\vdash A\preceq_S A}[NullProfileReflexive]
 \proofstepwidestar[1]{H_S}{\rA}
 \proofstepwidestar[2]{A\in K_S}{\rA}
 \proofstepwidestar[3]{X\in K_S}{\rA}
 \proofstepwidestar[4]{AX=z_S}{\rA}
 \proofstepwidestar[]{AX=z_S\to AX=z_S}{\rRI{4,4}}
 \proofstepwidestar[6]{XA=z_S}{\rA}
 \proofstepwidestar[]{XA=z_S\to XA=z_S}{\rRI{6,6}}
 \proofstepwidestar[]{\forall X\in K_S\;((AX=z_S\to AX=z_S)\land(XA=z_S\to XA=z_S))}{\rUI{\rRI{3,\rAI{5,7}}}}
 \proofstepwidestar[1,2]{A\preceq_S A}{\FormulaRefAuto{NullProfilePreorderDef}[def]{1,2,8}}
\closeproofpartwideindR
\Needspace{8\baselineskip}
\proofpartwideindR[Transitivität]{H_S\dsep A,B,C\in K_S\dsep A\preceq_S B\dsep B\preceq_S C\vdash A\preceq_S C}{H_S\dsep A,B,C\in K_S\dsep A\preceq_S B\dsep B\preceq_S C\vdash A\preceq_S C}[NullProfileTransitive]
 \proofstepwidestar[1]{H_S}{\rA}
 \proofstepwidestar[2]{A,B,C\in K_S}{\rA}
 \proofstepwidestar[3]{A\preceq_S B}{\rA}
 \proofstepwidestar[4]{B\preceq_S C}{\rA}
 \proofstepwidestar[5]{X\in K_S}{\rA}
 \proofstepwidestar[1,3,5]{(BX=z_S\to AX=z_S)\land(XB=z_S\to XA=z_S)}{\FormulaRefAuto{NullProfilePreorderDef}[def]{1,3,5}}
 \proofstepwidestar[1,4,5]{(CX=z_S\to BX=z_S)\land(XC=z_S\to XB=z_S)}{\FormulaRefAuto{NullProfilePreorderDef}[def]{1,4,5}}
 \proofstepwidestar[8]{CX=z_S}{\rA}
 \proofstepwidestar[1,3,4,5,8]{AX=z_S}{\rRE{\rAEa{6},\rRE{\rAEa{7},8}}}
 \proofstepwidestar[1,3,4,5]{CX=z_S\to AX=z_S}{\rRI{8,9}}
 \proofstepwidestar[11]{XC=z_S}{\rA}
 \proofstepwidestar[1,3,4,5,11]{XA=z_S}{\rRE{\rAEb{6},\rRE{\rAEb{7},11}}}
 \proofstepwidestar[1,3,4,5]{XC=z_S\to XA=z_S}{\rRI{11,12}}
 \proofstepwidestar[1,3,4]{\forall X\in K_S\;((CX=z_S\to AX=z_S)\land(XC=z_S\to XA=z_S))}{\rUI{\rRI{5,\rAI{10,13}}}}
 \proofstepwidestar[1,2,3,4]{A\preceq_S C}{\FormulaRefAuto{NullProfilePreorderDef}[def]{1,2,14}}
\closeproofpartwideindR
\Needspace{5\baselineskip}
\proofpartwide{Präordnung}
 \proofstepwidestar[1]{H_S}{\rA}
 \proofstepwidestar[2]{A\in K_S}{\rA}
 \proofstepwidestar[1,2]{A\preceq_S A}{\FormulaRefAuto{NullProfileReflexive}[thm]{1,2}}
 \proofstepwidestar[4]{B\in K_S}{\rA}
 \proofstepwidestar[5]{C\in K_S}{\rA}
 \proofstepwidestar[6]{A\preceq_S B\land B\preceq_S C}{\rA}
 \proofstepwidestar[1,2,4,5,6]{A\preceq_S C}{\FormulaRefAuto{NullProfileTransitive}[thm]{1,2,4,5,6}}
 \proofstepwidestar[1]{\operatorname{PrO}(K_S,\preceq_S)}{\FormulaRefAuto{PreorderDef}[def]{\rUI{\rRI{2,3}},\rUI{\rRI{2,\rUI{\rRI{4,\rUI{\rRI{5,\rRI{6,7}}}}}}}}}
\closeproofpartwide
\end{tabproofsplitwide}

\Needspace{9\baselineskip}
\FormulaThmDeltaK[Die Profilrelation ist eine Äquivalenzrelation]{%
 H_S\vdash\EqRel{K_S,\sim_S}
}{NullProfileEquivalence}{\DeltaRow{Mengen}{S\dsep0_S}}
\begin{tabproof}
 \proofstep{1}{H_S}{\rA}
 \proofstep{1}{\operatorname{PrO}(K_S,\preceq_S)}{\FormulaRefAuto{NullProfilePreorder}[thm]{1}}
 \proofstep{1}{\EqRel{K_S,\sim_S}}{\FormulaRefAuto{PreorderKernelEquivalence}[thm]{2}}
\end{tabproof}

\Needspace{14\baselineskip}
\FormulaDefDeltaK[Profilquotient und Singleton-Profilfasern]{%
 \begin{aligned}[t]
 Q_S&\coloneqq\QuotientSet{K_S}{\sim_S},&
 [A]_S&\coloneqq\EqClass{A}{\sim_S}{K_S},\\[-2pt]
 L_p&\coloneqq\{A\in K_S\mid[A]_S\le_Sp\},&
 D_p&\coloneqq\{s\in S\mid[\{s\}]_S\le_Sp\},\\[-2pt]
 C_p&\coloneqq\{s\in S\mid[\{s\}]_S=p\}.&&
 \end{aligned}
}{NullProfileSetsDef}{%
 \DeltaPrem{Halbgruppe mit Null}{H_S}
 \DeltaRow{Mengen}{S\dsep A\dsep p\dsep Q_S\dsep L_p\dsep D_p\dsep C_p}
 \DeltaRow{Quotientenordnung}{\le_S\ \text{aus }\FormulaRefAuto{PreorderQuotientOrderDef}[def]}}

\Needspace{9\baselineskip}
\FormulaThmDeltaK[Die Menge der Profile ist partiell geordnet]{%
 H_S\vdash\PartOrd{Q_S,\le_S}
}{NullProfilePartialOrder}{\DeltaRow{Mengen}{S\dsep Q_S}}
\begin{tabproof}
 \proofstep{1}{H_S}{\rA}
 \proofstep{1}{\operatorname{PrO}(K_S,\preceq_S)}{\FormulaRefAuto{NullProfilePreorder}[thm]{1}}
 \proofstep{1}{\PartOrd{Q_S,\le_S}}{\FormulaRefAuto{PreorderQuotientPartialOrder}[thm]{2}}
\end{tabproof}

\Needspace{9\baselineskip}
\FormulaThmDeltaK[Vergleich der Profile von Teilmengen]{%
 H_S\dsep A,B\in K_S\vdash([A]_S\le_S[B]_S\leftrightarrow A\preceq_S B)
}{NullProfileComparison}{\DeltaRow{Mengen}{S\dsep A\dsep B}}
\begin{tabproof}
 \proofstep{1}{H_S}{\rA}
 \proofstep{2}{A,B\in K_S}{\rA}
 \proofstep{1}{\operatorname{PrO}(K_S,\preceq_S)}{\FormulaRefAuto{NullProfilePreorder}[thm]{1}}
 \proofstep{1,2}{[A]_S\le_S[B]_S\leftrightarrow A\preceq_S B}{\FormulaRefAuto{PreorderQuotientComparison}[thm]{3,2}}
\end{tabproof}

\Needspace{9\baselineskip}
\FormulaThmDeltaK[Ein Profilvergleich zerfällt in Singletonvergleiche]{%
 H_S\dsep A,B\in K_S\vdash
 (A\preceq_S B\leftrightarrow\forall a\in A\;\{a\}\preceq_S B)
}{NullProfileSingletonCriterion}{\DeltaRow{Mengen}{S\dsep A\dsep B\dsep X\dsep a}}
\Needspace{6\baselineskip}
\begin{tabproofsplitwide}
\Needspace{5\baselineskip}
\proofpartwide{Von der Menge zu ihren Elementen}
 \proofstepwidestar[1]{H_S}{\rA}
 \proofstepwidestar[2]{A,B\in K_S}{\rA}
 \proofstepwidestar[1]{\SemigroupAlg{S}{\star}}{\FormulaRefAuto{SemigroupWithZeroBaseAxiom}[ax]{1}}
 \proofstepwidestar[4]{A\preceq_S B}{\rA}
 \proofstepwidestar[5]{a\in A}{\rA}
 \proofstepwidestar[6]{X\in K_S}{\rA}
 \proofstepwidestar[1,4,6]{(BX=z_S\to AX=z_S)\land(XB=z_S\to XA=z_S)}{\FormulaRefAuto{NullProfilePreorderDef}[def]{1,4,6}}
 \proofstepwidestar[8]{BX=z_S}{\rA}
 \proofstepwidestar[1,4,6,8]{AX=z_S}{\rRE{\rAEa{7},8}}
 \proofstepwidestar[1,2,4,5,6,8]{\{a\}X=z_S}{\FormulaRefAuto{PowerSemigroupLeftAtomization}[thm]{3,2,6,9,5}}
 \proofstepwidestar[1,2,4,5,6]{BX=z_S\to\{a\}X=z_S}{\rRI{8,10}}
 \proofstepwidestar[12]{XB=z_S}{\rA}
 \proofstepwidestar[1,4,6,12]{XA=z_S}{\rRE{\rAEb{7},12}}
 \proofstepwidestar[1,2,4,5,6,12]{X\{a\}=z_S}{\FormulaRefAuto{PowerSemigroupRightAtomization}[thm]{3,2,6,13,5}}
 \proofstepwidestar[1,2,4,5,6]{XB=z_S\to X\{a\}=z_S}{\rRI{12,14}}
 \proofstepwidestar[1,2,4,5]{\{a\}\preceq_S B}{\FormulaRefAuto{NullProfilePreorderDef}[def]{1,\rUI{\rRI{6,\rAI{11,15}}}}}
 \proofstepwidestar[1,2,4]{\forall a\in A\;\{a\}\preceq_S B}{\rUI{\rRI{5,16}}}
 \proofstepwidestar[1,2]{A\preceq_S B\to\forall a\in A\;\{a\}\preceq_S B}{\rRI{4,17}}
\closeproofpartwide
\Needspace{5\baselineskip}
\proofpartwide{Von den Elementen zur Menge}
 \setcounter{proofstepnr}{18}
 \proofstepwidestar[19]{\forall a\in A\;\{a\}\preceq_S B}{\rA}
 \proofstepwidestar[20]{X\in K_S}{\rA}
 \proofstepwidestar[21]{BX=z_S}{\rA}
 \proofstepwidestar[22]{a\in A}{\rA}
 \proofstepwidestar[19,22]{\{a\}\preceq_S B}{\rRE{\rUE{19},22}}
 \proofstepwidestar[1,19,20,22]{(BX=z_S\to\{a\}X=z_S)\land(XB=z_S\to X\{a\}=z_S)}{\FormulaRefAuto{NullProfilePreorderDef}[def]{1,23,20}}
 \proofstepwidestar[1,19,20,21,22]{\{a\}X=z_S}{\rRE{\rAEa{24},21}}
 \proofstepwidestar[1,19,20,21]{\forall a\in A\;\{a\}X=z_S}{\rUI{\rRI{22,25}}}
 \proofstepwidestar[1,2,19,20,21]{AX=z_S}{\FormulaRefAuto{PowerSemigroupLeftAtomization}[thm]{3,2,20,26}}
 \proofstepwidestar[1,2,19,20]{BX=z_S\to AX=z_S}{\rRI{21,27}}
 \proofstepwidestar[29]{XB=z_S}{\rA}
 \proofstepwidestar[1,19,20,22,29]{X\{a\}=z_S}{\rRE{\rAEb{24},29}}
 \proofstepwidestar[1,19,20,29]{\forall a\in A\;X\{a\}=z_S}{\rUI{\rRI{22,30}}}
 \proofstepwidestar[1,2,19,20,29]{XA=z_S}{\FormulaRefAuto{PowerSemigroupRightAtomization}[thm]{3,2,20,31}}
 \proofstepwidestar[1,2,19,20]{XB=z_S\to XA=z_S}{\rRI{29,32}}
 \proofstepwidestar[1,2,19]{A\preceq_S B}{\FormulaRefAuto{NullProfilePreorderDef}[def]{1,\rUI{\rRI{20,\rAI{28,33}}}}}
 \proofstepwidestar[1,2]{(\forall a\in A\;\{a\}\preceq_S B)\to A\preceq_S B}{\rRI{19,34}}
 \proofstepwidestar[1,2]{A\preceq_S B\leftrightarrow\forall a\in A\;\{a\}\preceq_S B}{\rLRI{18,35}}
\closeproofpartwide
\end{tabproofsplitwide}

\Needspace{9\baselineskip}
\FormulaThmDeltaK[Nulltests hängen nur von den beiden Profilen ab]{%
 \begin{aligned}[t]
 &H_S\dsep A,A',B,B'\in K_S\dsep A\sim_S A'\dsep B\sim_S B'\vdash{}\\[-2pt]
 &AB=z_S\leftrightarrow A'B'=z_S
 \end{aligned}
}{NullProfileProductInvariant}{\DeltaRow{Mengen}{S\dsep A\dsep A'\dsep B\dsep B'}}
\begin{tabproofwide}
 \proofstepwidestar[1]{H_S}{\rA}
 \proofstepwidestar[2]{A,A',B,B'\in K_S}{\rA}
 \proofstepwidestar[3]{A\sim_S A'}{\rA}
 \proofstepwidestar[4]{B\sim_S B'}{\rA}
 \proofstepwidestar[1,3]{A\preceq_S A'\land A'\preceq_S A}{\FormulaRefAuto{NullProfileEquivalenceDef}[def]{1,3}}
 \proofstepwidestar[1,4]{B\preceq_S B'\land B'\preceq_S B}{\FormulaRefAuto{NullProfileEquivalenceDef}[def]{1,4}}
 \proofstepwidestar[7]{AB=z_S}{\rA}
 \proofstepwidestar[1,2,3,7]{A'B=z_S}{\FormulaRefAuto{NullProfilePreorderDef}[def]{1,\rAEb{5},2,7}}
 \proofstepwidestar[1,2,3,4,7]{A'B'=z_S}{\FormulaRefAuto{NullProfilePreorderDef}[def]{1,\rAEb{6},2,8}}
 \proofstepwidestar[1,2,3,4]{AB=z_S\to A'B'=z_S}{\rRI{7,9}}
 \proofstepwidestar[11]{A'B'=z_S}{\rA}
 \proofstepwidestar[1,2,3,11]{AB'=z_S}{\FormulaRefAuto{NullProfilePreorderDef}[def]{1,\rAEa{5},2,11}}
 \proofstepwidestar[1,2,3,4,11]{AB=z_S}{\FormulaRefAuto{NullProfilePreorderDef}[def]{1,\rAEa{6},2,12}}
 \proofstepwidestar[1,2,3,4]{A'B'=z_S\to AB=z_S}{\rRI{11,13}}
 \proofstepwidestar[1,2,3,4]{AB=z_S\leftrightarrow A'B'=z_S}{\rLRI{10,14}}
\end{tabproofwide}

\Needspace{9\baselineskip}
\FormulaThmDeltaK[Das Hinzufügen der Null ändert ein Singletonprofil nicht]{%
 H_S\dsep c\in S\vdash\{c\}\sim_S\{0_S,c\}
}{NullProfileAdjoinZero}{\DeltaRow{Mengen}{S\dsep c\dsep x\dsep X}}
\Needspace{6\baselineskip}
\begin{tabproofsplitwide}
\Needspace{5\baselineskip}
\proofpartwide{Die beiden Nulltests}
 \proofstepwidestar[1]{H_S}{\rA}
 \proofstepwidestar[2]{c\in S}{\rA}
 \proofstepwidestar[1]{0_S\in S}{\FormulaRefAuto{SemigroupWithZeroCarrierAxiom}[ax]{1}}
 \proofstepwidestar[1,2]{\{0_S,c\}\subseteq S}{\FormulaRefAuto{a\in A,b\in A\vdash\{a,b\}\subseteq A}[thm]{3,2}}
 \proofstepwidestar[]{c\in\{0_S,c\}}{\FormulaRefAuto{b\in\{a,b\}}[thm]}
 \proofstepwidestar[]{\{0_S,c\}\neq\varnothing}{\FormulaRefAuto{a\in A\vdash A\neq\varnothing}[thm]{5}}
 \proofstepwidestar[1,2]{\{0_S,c\}\in K_S}{\FormulaRefAuto{NonemptyPowersetMembership}[thm]{\rAI{4,6}}}
 \proofstepwidestar[2]{\{c\}\in K_S}{\FormulaRefAuto{SingletonPowerMember}[thm]{2}}
 \proofstepwidestar[1]{\SemigroupAlg{S}{\star}}{\FormulaRefAuto{SemigroupWithZeroBaseAxiom}[ax]{1}}
 \proofstepwidestar[10]{X\in K_S}{\rA}
 \proofstepwidestar[1]{\SemigroupWithZeroAlg{K_S}{\star_{\mathcal P}}{z_S}}{\FormulaRefAuto{PowerSemigroupWithZero}[thm]{1}}
 \proofstepwidestar[1,10]{z_SX=z_S\land Xz_S=z_S}{\FormulaRefAuto{SemigroupWithZeroDef}[def]{11,10}}
 \proofstepwidestar[13]{\{c\}X=z_S}{\rA}
 \proofstepwidestar[14]{x\in\{0_S,c\}}{\rA}
 \proofstepwidestar[14]{x=0_S\lor x=c}{\FormulaRefAuto{x\in\{a,b\}\eqvdash(x=a\lor x=b)}[thm]{14}}
 \proofstepwidestar[16]{x=0_S}{\rA}
 \proofstepwidestar[1,10,16]{\{x\}X=z_S}{\rIE{\FormulaRefAuto{a=b\vdash b=a}[thm]{16},\rAEa{12}}}
 \proofstepwidestar[18]{x=c}{\rA}
 \proofstepwidestar[13,18]{\{x\}X=z_S}{\rIE{\FormulaRefAuto{a=b\vdash b=a}[thm]{18},13}}
 \proofstepwidestar[1,10,13,14]{\{x\}X=z_S}{\rOE{15,[16]\vdots17,[18]\vdots19}}
 \proofstepwidestar[1,10,13]{\forall x\in\{0_S,c\}\;\{x\}X=z_S}{\rUI{\rRI{14,20}}}
 \proofstepwidestar[1,2,10,13]{\{0_S,c\}X=z_S}{\FormulaRefAuto{PowerSemigroupLeftAtomization}[thm]{9,7,10,21}}
 \proofstepwidestar[1,2,10]{\{c\}X=z_S\to\{0_S,c\}X=z_S}{\rRI{13,22}}
 \proofstepwidestar[24]{X\{c\}=z_S}{\rA}
 \proofstepwidestar[1,10,16]{X\{x\}=z_S}{\rIE{\FormulaRefAuto{a=b\vdash b=a}[thm]{16},\rAEb{12}}}
 \proofstepwidestar[18,24]{X\{x\}=z_S}{\rIE{\FormulaRefAuto{a=b\vdash b=a}[thm]{18},24}}
 \proofstepwidestar[1,10,14,24]{X\{x\}=z_S}{\rOE{15,[16]\vdots25,[18]\vdots26}}
 \proofstepwidestar[1,10,24]{\forall x\in\{0_S,c\}\;X\{x\}=z_S}{\rUI{\rRI{14,27}}}
 \proofstepwidestar[1,2,10,24]{X\{0_S,c\}=z_S}{\FormulaRefAuto{PowerSemigroupRightAtomization}[thm]{9,7,10,28}}
 \proofstepwidestar[1,2,10]{X\{c\}=z_S\to X\{0_S,c\}=z_S}{\rRI{24,29}}
\closeproofpartwide
\Needspace{5\baselineskip}
\proofpartwide{Beide Profilrichtungen}
 \setcounter{proofstepnr}{30}
 \proofstepwidestar[1,2]{\{0_S,c\}\preceq_S\{c\}}{\FormulaRefAuto{NullProfilePreorderDef}[def]{1,\rUI{\rRI{10,\rAI{23,30}}}}}
 \proofstepwidestar[32]{\{0_S,c\}X=z_S}{\rA}
 \proofstepwidestar[1,2,10,32]{\{c\}X=z_S}{\FormulaRefAuto{PowerSemigroupLeftAtomization}[thm]{9,7,10,32,5}}
 \proofstepwidestar[1,2,10]{\{0_S,c\}X=z_S\to\{c\}X=z_S}{\rRI{32,33}}
 \proofstepwidestar[35]{X\{0_S,c\}=z_S}{\rA}
 \proofstepwidestar[1,2,10,35]{X\{c\}=z_S}{\FormulaRefAuto{PowerSemigroupRightAtomization}[thm]{9,7,10,35,5}}
 \proofstepwidestar[1,2,10]{X\{0_S,c\}=z_S\to X\{c\}=z_S}{\rRI{35,36}}
 \proofstepwidestar[1,2]{\{c\}\preceq_S\{0_S,c\}}{\FormulaRefAuto{NullProfilePreorderDef}[def]{1,\rUI{\rRI{10,\rAI{34,37}}}}}
 \proofstepwidestar[1,2]{\{c\}\sim_S\{0_S,c\}}{\FormulaRefAuto{NullProfileEquivalenceDef}[def]{1,\rAI{38,31}}}
\closeproofpartwide
\end{tabproofsplitwide}

\Needspace{16\baselineskip}
\section{Untere Mengen und Grundelemente}

\Needspace{9\baselineskip}
\FormulaThmDeltaK[Vergleich mit einem beliebigen Profil]{%
 H_S\dsep p\in Q_S\dsep A\in K_S\vdash
 ([A]_S\le_Sp\leftrightarrow\forall a\in A\;[\{a\}]_S\le_Sp)
}{NullProfileLowerSingletonCriterion}{\DeltaRow{Mengen}{S\dsep p\dsep A\dsep B\dsep a}}
\begin{tabproofwide}
 \proofstepwidestar[1]{H_S}{\rA}
 \proofstepwidestar[2]{p\in Q_S}{\rA}
 \proofstepwidestar[3]{A\in K_S}{\rA}
 \proofstepwidestar[1]{\EqRel{K_S,\sim_S}}{\FormulaRefAuto{NullProfileEquivalence}[thm]{1}}
 \proofstepwidestar[1,2]{\exists B\in K_S\;p=[B]_S}{\FormulaRefAuto{QuotientRepresentative}[thm]{4,2}}
 \proofstepwidestar[6]{B\in K_S\land p=[B]_S}{\rA}
 \proofstepwidestar[1,3,6]{[A]_S\le_S[B]_S\leftrightarrow A\preceq_S B}{\FormulaRefAuto{NullProfileComparison}[thm]{1,3,6}}
 \proofstepwidestar[1,3,6]{A\preceq_S B\leftrightarrow\forall a\in A\;\{a\}\preceq_S B}{\FormulaRefAuto{NullProfileSingletonCriterion}[thm]{1,3,6}}
 \proofstepwidestar[9]{a\in A}{\rA}
 \proofstepwidestar[3,9]{\{a\}\in K_S}{\FormulaRefAuto{SingletonPowerMember}[thm]{\FormulaRefAuto{PowerSemigroupCarrierElement}[thm]{3,9}}}
 \proofstepwidestar[1,3,6,9]{\{a\}\preceq_S B\leftrightarrow[\{a\}]_S\le_S[B]_S}{\FormulaRefAuto{P\leftrightarrow Q\vdash Q\leftrightarrow P}[thm]{\FormulaRefAuto{NullProfileComparison}[thm]{1,10,6}}}
 \proofstepwidestar[1,3,6]{\forall a\in A\;(\{a\}\preceq_S B\leftrightarrow[\{a\}]_S\le_S[B]_S)}{\rUI{\rRI{9,11}}}
 \proofstepwidestar[13]{\forall a\in A\;\{a\}\preceq_S B}{\rA}
 \proofstepwidestar[9,13]{\{a\}\preceq_S B}{\rRE{\rUE{13},9}}
 \proofstepwidestar[1,3,6,9,13]{[\{a\}]_S\le_S[B]_S}{\FormulaRefAuto{P\leftrightarrow Q,P\vdash Q}[thm]{11,14}}
 \proofstepwidestar[1,3,6,13]{\forall a\in A\;[\{a\}]_S\le_S[B]_S}{\rUI{\rRI{9,15}}}
 \proofstepwidestar[1,3,6]{(\forall a\in A\;\{a\}\preceq_S B)\to\forall a\in A\;[\{a\}]_S\le_S[B]_S}{\rRI{13,16}}
 \proofstepwidestar[18]{\forall a\in A\;[\{a\}]_S\le_S[B]_S}{\rA}
 \proofstepwidestar[9,18]{[\{a\}]_S\le_S[B]_S}{\rRE{\rUE{18},9}}
 \proofstepwidestar[1,3,6,9,18]{\{a\}\preceq_S B}{\FormulaRefAuto{P\leftrightarrow Q,Q\vdash P}[thm]{11,19}}
 \proofstepwidestar[1,3,6,18]{\forall a\in A\;\{a\}\preceq_S B}{\rUI{\rRI{9,20}}}
 \proofstepwidestar[1,3,6]{(\forall a\in A\;[\{a\}]_S\le_S[B]_S)\to\forall a\in A\;\{a\}\preceq_S B}{\rRI{18,21}}
 \proofstepwidestar[1,3,6]{(\forall a\in A\;\{a\}\preceq_S B)\leftrightarrow\forall a\in A\;[\{a\}]_S\le_S[B]_S}{\rLRI{17,22}}
 \proofstepwidestar[1,3,6]{[A]_S\le_S[B]_S\leftrightarrow\forall a\in A\;[\{a\}]_S\le_S[B]_S}{\FormulaRefAuto{P\leftrightarrow Q,Q\leftrightarrow R\vdash P\leftrightarrow R}[thm]{\FormulaRefAuto{P\leftrightarrow Q,Q\leftrightarrow R\vdash P\leftrightarrow R}[thm]{7,8},23}}
 \proofstepwidestar[1,3,6]{[A]_S\le_Sp\leftrightarrow\forall a\in A\;[\{a\}]_S\le_Sp}{\rIE{\FormulaRefAuto{a=b\vdash b=a}[thm]{\rAEb{6}},24}}
 \proofstepwidestar[1,2,3]{[A]_S\le_Sp\leftrightarrow\forall a\in A\;[\{a\}]_S\le_Sp}{\rEE{5,6,25}}
\end{tabproofwide}

\Needspace{9\baselineskip}
\FormulaThmDeltaK[Untere Mengen sind nichtleere Potenzmengen]{%
 H_S\dsep p\in Q_S\vdash L_p=\PowNE{D_p}
}{NullProfileLowerPowerSet}{\DeltaRow{Mengen}{S\dsep p\dsep A\dsep a\dsep L_p\dsep D_p}}
\Needspace{6\baselineskip}
\begin{tabproofsplitwide}
\Needspace{5\baselineskip}
\proofpartwide{Von der unteren Menge zur Potenzmenge}
 \proofstepwidestar[1]{H_S}{\rA}
 \proofstepwidestar[2]{p\in Q_S}{\rA}
 \proofstepwidestar[3]{A\in L_p}{\rA}
 \proofstepwidestar[1,3]{A\in K_S\land[A]_S\le_Sp}{\FormulaRefAuto{NullProfileSetsDef}[def]{1,3}}
 \proofstepwidestar[1,2,3]{\forall a\in A\;[\{a\}]_S\le_Sp}{\FormulaRefAuto{NullProfileLowerSingletonCriterion}[thm]{1,2,\rAEa{4},\rAEb{4}}}
 \proofstepwidestar[6]{a\in A}{\rA}
 \proofstepwidestar[1,3,6]{a\in S}{\FormulaRefAuto{PowerSemigroupCarrierElement}[thm]{\rAEa{4},6}}
 \proofstepwidestar[1,2,3,6]{[\{a\}]_S\le_Sp}{\rRE{\rUE{5},6}}
 \proofstepwidestar[1,2,3,6]{a\in D_p}{\FormulaRefAuto{NullProfileSetsDef}[def]{1,\rAI{7,8}}}
 \proofstepwidestar[1,2,3]{A\subseteq D_p}{\FormulaRefAuto{SubsetIntroductionRule}[thm]{[6]\vdots9}}
 \proofstepwidestar[1,3]{A\neq\varnothing}{\rAEb{\FormulaRefAuto{NonemptyPowersetMembership}[thm]{\rAEa{4}}}}
 \proofstepwidestar[1,2,3]{A\in\PowNE{D_p}}{\FormulaRefAuto{NonemptyPowersetMembership}[thm]{\rAI{10,11}}}
 \proofstepwidestar[1,2]{L_p\subseteq\PowNE{D_p}}{\FormulaRefAuto{SubsetIntroductionRule}[thm]{[3]\vdots12}}
\closeproofpartwide
\Needspace{5\baselineskip}
\proofpartwide{Von der Potenzmenge zur unteren Menge}
 \setcounter{proofstepnr}{13}
 \proofstepwidestar[14]{A\in\PowNE{D_p}}{\rA}
 \proofstepwidestar[14]{A\subseteq D_p\land A\neq\varnothing}{\FormulaRefAuto{NonemptyPowersetMembership}[thm]{14}}
 \proofstepwidestar[1]{D_p\subseteq S}{\FormulaRefAuto{\{x\in A\mid P(x)\}\subseteq A}[thm]{\FormulaRefAuto{NullProfileSetsDef}[def]{1}}}
 \proofstepwidestar[1,14]{A\subseteq S}{\FormulaRefAuto{A\subseteq B,B\subseteq C\vdash A\subseteq C}[thm]{\rAEa{15},16}}
 \proofstepwidestar[1,14]{A\in K_S}{\FormulaRefAuto{NonemptyPowersetMembership}[thm]{\rAI{17,\rAEb{15}}}}
 \proofstepwidestar[19]{a\in A}{\rA}
 \proofstepwidestar[14,19]{a\in D_p}{\FormulaRefAuto{A\subseteq B\dsep x\in A\vdash x\in B}[thm]{\rAEa{15},19}}
 \proofstepwidestar[1,14,19]{[\{a\}]_S\le_Sp}{\rAEb{\FormulaRefAuto{NullProfileSetsDef}[def]{1,20}}}
 \proofstepwidestar[1,14]{\forall a\in A\;[\{a\}]_S\le_Sp}{\rUI{\rRI{19,21}}}
 \proofstepwidestar[1,2,14]{[A]_S\le_Sp}{\FormulaRefAuto{NullProfileLowerSingletonCriterion}[thm]{1,2,18,22}}
 \proofstepwidestar[1,2,14]{A\in L_p}{\FormulaRefAuto{NullProfileSetsDef}[def]{1,\rAI{18,23}}}
 \proofstepwidestar[1,2]{\PowNE{D_p}\subseteq L_p}{\FormulaRefAuto{SubsetIntroductionRule}[thm]{[14]\vdots24}}
 \proofstepwidestar[1,2]{L_p=\PowNE{D_p}}{\FormulaRefAuto{A\subseteq B,B\subseteq A\vdash A=B}[thm]{13,25}}
\closeproofpartwide
\end{tabproofsplitwide}

\Needspace{16\baselineskip}
\section{Transport durch Isomorphismen}

\Needspace{9\baselineskip}
\FormulaThmDeltaK[Transport der Nullproduktrelation]{%
 H_S\dsep H_T\dsep I\dsep A,B\in K_S\vdash
 (AB=z_S\leftrightarrow\Phi(A)\Phi(B)=z_T)
}{NullProfileIsoZeroProduct}{\DeltaRow{Mengen}{S\dsep T\dsep A\dsep B}}
\begin{tabproofwide}
 \proofstepwidestar[1]{H_S}{\rA}
 \proofstepwidestar[2]{H_T}{\rA}
 \proofstepwidestar[3]{I}{\rA}
 \proofstepwidestar[4]{A,B\in K_S}{\rA}
 \proofstepwidestar[1,2,3]{\Phi(z_S)=z_T}{\FormulaRefAuto{PowerSemigroupZeroIsoValue}[thm]{1,2,3}}
 \proofstepwidestar[3]{\SemigroupAlg{K_S}{\star_{\mathcal P}}}{\FormulaRefAuto{SemigroupIsoSource}[thm]{3}}
 \proofstepwidestar[3,4]{AB\in K_S}{\FormulaRefAuto{SemigroupClosureAxiom}[ax]{6,4}}
 \proofstepwidestar[1]{z_S\in K_S}{\FormulaRefAuto{SemigroupWithZeroCarrierAxiom}[ax]{\FormulaRefAuto{PowerSemigroupWithZero}[thm]{1}}}
 \proofstepwidestar[3,4]{\Phi(AB)=\Phi(A)\Phi(B)}{\FormulaRefAuto{SemigroupIsoOperation}[thm]{3,4}}
 \proofstepwidestar[1,3,4]{AB=z_S\leftrightarrow\Phi(AB)=\Phi(z_S)}{\FormulaRefAuto{SemigroupIsoEqualityReflection}[thm]{3,7,8}}
 \proofstepwidestar[1,2,3,4]{AB=z_S\leftrightarrow\Phi(A)\Phi(B)=z_T}{\rIE{5,\rIE{9,10}}}
\end{tabproofwide}

\Needspace{9\baselineskip}
\FormulaThmDeltaK[Isomorphismen erhalten den Profilvergleich]{%
 H_S\dsep H_T\dsep I\dsep A,B\in K_S\dsep A\preceq_S B\vdash
 \Phi(A)\preceq_T\Phi(B)
}{NullProfileIsoPreorderForward}{\DeltaRow{Mengen}{S\dsep T\dsep A\dsep B\dsep X\dsep Y}}
\Needspace{6\baselineskip}
\begin{tabproofsplitwide}
\Needspace{5\baselineskip}
\proofpartwide{Urbild eines beliebigen Nulltests}
 \proofstepwidestar[1]{H_S}{\rA}
 \proofstepwidestar[2]{H_T}{\rA}
 \proofstepwidestar[3]{I}{\rA}
 \proofstepwidestar[4]{A,B\in K_S}{\rA}
 \proofstepwidestar[5]{A\preceq_S B}{\rA}
 \proofstepwidestar[6]{Y\in K_T}{\rA}
 \proofstepwidestar[3]{\Phi:K_S\bij K_T}{\FormulaRefAuto{SemigroupIsoBijectiveAxiom}[ax]{3}}
 \proofstepwidestar[3]{\Phi^{-1}:K_T\to K_S}{\FormulaRefAuto{InverseFunctionType}[thm]{7}}
 \proofstepwidestar[3,6]{\Phi^{-1}(Y)\in K_S}{\FormulaRefAuto{F\colon A\to B\dsep x\in A\vdash F(x)\in B}[thm]{8,6}}
 \proofstepwidestar[3,6]{\Phi(\Phi^{-1}(Y))=Y}{\FormulaRefAuto{InverseFunctionRightInverse}[thm]{7,6}}
 \proofstepwidestar[1,3,5,6]{\begin{aligned}[t]&(B\Phi^{-1}(Y)=z_S\to A\Phi^{-1}(Y)=z_S)\\[-2pt]&\land(\Phi^{-1}(Y)B=z_S\to\Phi^{-1}(Y)A=z_S)\end{aligned}}{\FormulaRefAuto{NullProfilePreorderDef}[def]{1,5,9}}
\closeproofpartwide
\Needspace{5\baselineskip}
\proofpartwide{Rechter und linker Test}
 \setcounter{proofstepnr}{11}
 \proofstepwidestar[12]{\Phi(B)Y=z_T}{\rA}
 \proofstepwidestar[1,2,3,4,6,12]{B\Phi^{-1}(Y)=z_S}{\FormulaRefAuto{NullProfileIsoZeroProduct}[thm]{1,2,3,4,9,\rIE{\FormulaRefAuto{a=b\vdash b=a}[thm]{10},12}}}
 \proofstepwidestar[1,2,3,4,5,6,12]{A\Phi^{-1}(Y)=z_S}{\rRE{\rAEa{11},13}}
 \proofstepwidestar[1,2,3,4,5,6,12]{\Phi(A)Y=z_T}{\rIE{10,\FormulaRefAuto{NullProfileIsoZeroProduct}[thm]{1,2,3,4,9,14}}}
 \proofstepwidestar[1,2,3,4,5,6]{\Phi(B)Y=z_T\to\Phi(A)Y=z_T}{\rRI{12,15}}
 \proofstepwidestar[17]{Y\Phi(B)=z_T}{\rA}
 \proofstepwidestar[1,2,3,4,6,17]{\Phi^{-1}(Y)B=z_S}{\FormulaRefAuto{NullProfileIsoZeroProduct}[thm]{1,2,3,9,4,\rIE{\FormulaRefAuto{a=b\vdash b=a}[thm]{10},17}}}
 \proofstepwidestar[1,2,3,4,5,6,17]{\Phi^{-1}(Y)A=z_S}{\rRE{\rAEb{11},18}}
 \proofstepwidestar[1,2,3,4,5,6,17]{Y\Phi(A)=z_T}{\rIE{10,\FormulaRefAuto{NullProfileIsoZeroProduct}[thm]{1,2,3,9,4,19}}}
 \proofstepwidestar[1,2,3,4,5,6]{Y\Phi(B)=z_T\to Y\Phi(A)=z_T}{\rRI{17,20}}
 \proofstepwidestar[1,2,3,4,5]{\Phi(A)\preceq_T\Phi(B)}{\FormulaRefAuto{NullProfilePreorderDef}[def]{2,\rUI{\rRI{6,\rAI{16,21}}}}}
\closeproofpartwide
\end{tabproofsplitwide}

\Needspace{9\baselineskip}
\FormulaThmDeltaK[Isomorphismen erhalten und reflektieren Profilvergleiche]{%
 H_S\dsep H_T\dsep I\dsep A,B\in K_S\vdash
 (A\preceq_S B\leftrightarrow\Phi(A)\preceq_T\Phi(B))
}{NullProfileIsoPreorder}{\DeltaRow{Mengen}{S\dsep T\dsep A\dsep B}}
\begin{tabproofwide}
 \proofstepwidestar[1]{H_S}{\rA}
 \proofstepwidestar[2]{H_T}{\rA}
 \proofstepwidestar[3]{I}{\rA}
 \proofstepwidestar[4]{A,B\in K_S}{\rA}
 \proofstepwidestar[5]{A\preceq_S B}{\rA}
 \proofstepwidestar[1,2,3,4,5]{\Phi(A)\preceq_T\Phi(B)}{\FormulaRefAuto{NullProfileIsoPreorderForward}[thm]{1,2,3,4,5}}
 \proofstepwidestar[1,2,3,4]{A\preceq_S B\to\Phi(A)\preceq_T\Phi(B)}{\rRI{5,6}}
 \proofstepwidestar[3]{\SemigroupIso{\Phi^{-1}}{K_T,\diamond_{\mathcal P}}{K_S,\star_{\mathcal P}}}{\FormulaRefAuto{SemigroupIsoInverse}[thm]{3}}
 \proofstepwidestar[3]{\Phi:K_S\bij K_T}{\FormulaRefAuto{SemigroupIsoBijectiveAxiom}[ax]{3}}
 \proofstepwidestar[3,4]{\Phi(A),\Phi(B)\in K_T}{\FormulaRefAuto{F\colon A\to B\dsep x\in A\vdash F(x)\in B}[thm]{\FormulaRefAuto{Bijektive Funktion}[def]{9},4}}
 \proofstepwidestar[11]{\Phi(A)\preceq_T\Phi(B)}{\rA}
 \proofstepwidestar[1,2,3,4,11]{\Phi^{-1}(\Phi(A))\preceq_S\Phi^{-1}(\Phi(B))}{\FormulaRefAuto{NullProfileIsoPreorderForward}[thm]{2,1,8,10,11}}
 \proofstepwidestar[3,4]{\Phi^{-1}(\Phi(A))=A\land\Phi^{-1}(\Phi(B))=B}{\FormulaRefAuto{InverseFunctionLeftInverse}[thm]{9,4}}
 \proofstepwidestar[1,2,3,4,11]{A\preceq_S B}{\rIE{\rAEa{13},\rAEb{13},12}}
 \proofstepwidestar[1,2,3,4]{\Phi(A)\preceq_T\Phi(B)\to A\preceq_S B}{\rRI{11,14}}
 \proofstepwidestar[1,2,3,4]{A\preceq_S B\leftrightarrow\Phi(A)\preceq_T\Phi(B)}{\rLRI{7,15}}
\end{tabproofwide}

\Needspace{9\baselineskip}
\FormulaThmDeltaK[Transport der Profiläquivalenz]{%
 H_S\dsep H_T\dsep I\dsep A,B\in K_S\vdash
 (A\sim_S B\leftrightarrow\Phi(A)\sim_T\Phi(B))
}{NullProfileIsoEquivalence}{\DeltaRow{Mengen}{S\dsep T\dsep A\dsep B}}
\begin{tabproofwide}
 \proofstepwidestar[1]{H_S}{\rA}
 \proofstepwidestar[2]{H_T}{\rA}
 \proofstepwidestar[3]{I}{\rA}
 \proofstepwidestar[4]{A,B\in K_S}{\rA}
 \proofstepwidestar[1,2,3,4]{A\preceq_S B\leftrightarrow\Phi(A)\preceq_T\Phi(B)}{\FormulaRefAuto{NullProfileIsoPreorder}[thm]{1,2,3,4}}
 \proofstepwidestar[1,2,3,4]{B\preceq_S A\leftrightarrow\Phi(B)\preceq_T\Phi(A)}{\FormulaRefAuto{NullProfileIsoPreorder}[thm]{1,2,3,4}}
 \proofstepwidestar[1]{A\sim_S B\leftrightarrow(A\preceq_S B\land B\preceq_S A)}{\FormulaRefAuto{NullProfileEquivalenceDef}[def]{1}}
 \proofstepwidestar[2]{\Phi(A)\sim_T\Phi(B)\leftrightarrow(\Phi(A)\preceq_T\Phi(B)\land\Phi(B)\preceq_T\Phi(A))}{\FormulaRefAuto{NullProfileEquivalenceDef}[def]{2}}
 \proofstepwidestar[1,2,3,4]{A\sim_S B\leftrightarrow(\Phi(A)\preceq_T\Phi(B)\land\Phi(B)\preceq_T\Phi(A))}{\rLRS{6,\rLRS{5,7}}}
 \proofstepwidestar[1,2,3,4]{A\sim_S B\leftrightarrow\Phi(A)\sim_T\Phi(B)}{\FormulaRefAuto{P\leftrightarrow Q,R\leftrightarrow Q\vdash P\leftrightarrow R}[thm]{9,8}}
\end{tabproofwide}

Zur Konstruktion der Profilabbildung seien
\(\pi_S=\QuotProj{K_S}{\sim_S}\) und
\(\pi_T=\QuotProj{K_T}{\sim_T}\). Die folgende Konstruktion wählt
keine Repräsentanten aus den Profilklassen.

\Needspace{9\baselineskip}
\FormulaThmDeltaK[Existenz der induzierten Profilabbildung]{%
 H_S\dsep H_T\dsep I\vdash
 \exists!g\;(g:Q_S\to Q_T\land g\circ\pi_S=\pi_T\circ\Phi)
}{NullProfileInducedMapExistence}{\DeltaRow{Mengen}{S\dsep T\dsep A\dsep B}}
\begin{tabproofwide}
 \proofstepwidestar[1]{H_S}{\rA}
 \proofstepwidestar[2]{H_T}{\rA}
 \proofstepwidestar[3]{I}{\rA}
 \proofstepwidestar[1]{\EqRel{K_S,\sim_S}}{\FormulaRefAuto{NullProfileEquivalence}[thm]{1}}
 \proofstepwidestar[2]{\EqRel{K_T,\sim_T}}{\FormulaRefAuto{NullProfileEquivalence}[thm]{2}}
 \proofstepwidestar[2]{\pi_T:K_T\to Q_T}{\FormulaRefAuto{QuotProjFunction}[thm]{5}}
 \proofstepwidestar[3]{\Phi:K_S\to K_T}{\FormulaRefAuto{SemigroupIsoFunction}[thm]{3}}
 \proofstepwidestar[2,3]{\pi_T\circ\Phi:K_S\to Q_T}{\FormulaRefAuto{CompositionDef}[def]{7,6}}
 \proofstepwidestar[9]{A,B\in K_S}{\rA}
 \proofstepwidestar[10]{A\sim_S B}{\rA}
 \proofstepwidestar[1,2,3,9,10]{\Phi(A)\sim_T\Phi(B)}{\FormulaRefAuto{NullProfileIsoEquivalence}[thm]{1,2,3,9,10}}
 \proofstepwidestar[3,9]{\Phi(A),\Phi(B)\in K_T}{\FormulaRefAuto{F\colon A\to B\dsep x\in A\vdash F(x)\in B}[thm]{7,9}}
 \proofstepwidestar[1,2,3,9,10]{[\Phi(A)]_T=[\Phi(B)]_T}{\FormulaRefAuto{EqClassEqualFromRel}[thm]{5,12,11}}
 \proofstepwidestar[2,3,9]{\pi_T(\Phi(A))=[\Phi(A)]_T\land\pi_T(\Phi(B))=[\Phi(B)]_T}{\FormulaRefAuto{QuotProjValue}[thm]{5,12}}
 \proofstepwidestar[2,3,9]{(\pi_T\circ\Phi)(A)=\pi_T(\Phi(A))\land(\pi_T\circ\Phi)(B)=\pi_T(\Phi(B))}{\FormulaRefAuto{CompositionDef}[def]{7,6,9}}
 \proofstepwidestar[1,2,3,9,10]{(\pi_T\circ\Phi)(A)=(\pi_T\circ\Phi)(B)}{\rIE{\FormulaRefAuto{a=b\vdash b=a}[thm]{\rAEa{14}},\FormulaRefAuto{a=b\vdash b=a}[thm]{\rAEb{14}},\FormulaRefAuto{a=b\vdash b=a}[thm]{\rAEa{15}},\FormulaRefAuto{a=b\vdash b=a}[thm]{\rAEb{15}},13}}
 \proofstepwidestar[1,2,3]{\forall A,B\in K_S\;(A\sim_S B\to(\pi_T\circ\Phi)(A)=(\pi_T\circ\Phi)(B))}{\rUI{\rUI{\rRI{9,\rRI{10,16}}}}}
 \proofstepwidestar[1,2,3]{\exists!g\;(g:Q_S\to Q_T\land g\circ\pi_S=\pi_T\circ\Phi)}{\FormulaRefAuto{QuotientFunctionDescent}[thm]{4,8,17}}
\end{tabproofwide}

\Needspace{12\baselineskip}
\FormulaDefDeltaK[Die induzierte Profilabbildung]{%
 \psi\coloneqq\iota g\;(g:Q_S\to Q_T\land g\circ\pi_S=\pi_T\circ\Phi)
}{NullProfileInducedMapDef}{%
 \DeltaPrem{Quellhalbgruppe mit Null}{H_S}
 \DeltaPrem{Zielhalbgruppe mit Null}{H_T}
 \DeltaPrem{Potenzisomorphismus}{I}
 \DeltaRow{Eindeutigkeit}{\FormulaRefAuto{NullProfileInducedMapExistence}[thm]}}

\Needspace{9\baselineskip}
\FormulaThmDeltaK[Typ und Repräsentantenwert der Profilabbildung]{%
 H_S\dsep H_T\dsep I\vdash
 \psi:Q_S\to Q_T\land\forall A\in K_S\;(\psi([A]_S)=[\Phi(A)]_T)
}{NullProfileInducedMapValue}{\DeltaRow{Mengen}{S\dsep T\dsep A}}
\begin{tabproofwide}
 \proofstepwidestar[1]{H_S}{\rA}
 \proofstepwidestar[2]{H_T}{\rA}
 \proofstepwidestar[3]{I}{\rA}
 \proofstepwidestar[1,2,3]{\psi:Q_S\to Q_T\land\psi\circ\pi_S=\pi_T\circ\Phi}{\FormulaRefAuto{NullProfileInducedMapDef}[def]{\FormulaRefAuto{NullProfileInducedMapExistence}[thm]{1,2,3}}}
 \proofstepwidestar[5]{A\in K_S}{\rA}
 \proofstepwidestar[1]{\EqRel{K_S,\sim_S}}{\FormulaRefAuto{NullProfileEquivalence}[thm]{1}}
 \proofstepwidestar[2]{\EqRel{K_T,\sim_T}}{\FormulaRefAuto{NullProfileEquivalence}[thm]{2}}
 \proofstepwidestar[3]{\Phi:K_S\to K_T}{\FormulaRefAuto{SemigroupIsoFunction}[thm]{3}}
 \proofstepwidestar[3,5]{\Phi(A)\in K_T}{\FormulaRefAuto{F\colon A\to B\dsep x\in A\vdash F(x)\in B}[thm]{8,5}}
 \proofstepwidestar[1]{\pi_S:K_S\to Q_S}{\FormulaRefAuto{QuotProjFunction}[thm]{6}}
 \proofstepwidestar[2]{\pi_T:K_T\to Q_T}{\FormulaRefAuto{QuotProjFunction}[thm]{7}}
 \proofstepwidestar[1,5]{\pi_S(A)=[A]_S}{\FormulaRefAuto{QuotProjValue}[thm]{6,5}}
 \proofstepwidestar[2,3,5]{\pi_T(\Phi(A))=[\Phi(A)]_T}{\FormulaRefAuto{QuotProjValue}[thm]{7,9}}
 \proofstepwidestar[1,2,3,5]{(\psi\circ\pi_S)(A)=\psi(\pi_S(A))}{\FormulaRefAuto{CompositionDef}[def]{10,\rAEa{4},5}}
 \proofstepwidestar[2,3,5]{(\pi_T\circ\Phi)(A)=\pi_T(\Phi(A))}{\FormulaRefAuto{CompositionDef}[def]{8,11,5}}
 \proofstepwidestar[1,2,3,5]{\psi([A]_S)=[\Phi(A)]_T}{\FormulaRefAuto{a=b\vdash b=a}[thm]{\rIE{\rAEb{4},15,12,13,14}}}
 \proofstepwidestar[1,2,3]{\forall A\in K_S\;(\psi([A]_S)=[\Phi(A)]_T)}{\rUI{\rRI{5,16}}}
 \proofstepwidestar[1,2,3]{\psi:Q_S\to Q_T\land\forall A\in K_S\;(\psi([A]_S)=[\Phi(A)]_T)}{\rAI{\rAEa{4},17}}
\end{tabproofwide}

\Needspace{9\baselineskip}
\FormulaThmDeltaK[Die Profilabbildung erhält und reflektiert die Ordnung]{%
 H_S\dsep H_T\dsep I\dsep p,q\in Q_S\vdash
 (p\le_Sq\leftrightarrow\psi(p)\le_T\psi(q))
}{NullProfileIsoOrder}{\DeltaRow{Mengen}{S\dsep T\dsep p\dsep q\dsep A\dsep B}}
\begin{tabproofwide}
 \proofstepwidestar[1]{H_S}{\rA}
 \proofstepwidestar[2]{H_T}{\rA}
 \proofstepwidestar[3]{I}{\rA}
 \proofstepwidestar[4]{p,q\in Q_S}{\rA}
 \proofstepwidestar[1]{\EqRel{K_S,\sim_S}}{\FormulaRefAuto{NullProfileEquivalence}[thm]{1}}
 \proofstepwidestar[1,4]{\exists A\in K_S\;p=[A]_S}{\FormulaRefAuto{QuotientRepresentative}[thm]{5,4}}
 \proofstepwidestar[1,4]{\exists B\in K_S\;q=[B]_S}{\FormulaRefAuto{QuotientRepresentative}[thm]{5,4}}
 \proofstepwidestar[8]{A\in K_S\land p=[A]_S}{\rA}
 \proofstepwidestar[9]{B\in K_S\land q=[B]_S}{\rA}
 \proofstepwidestar[1,8,9]{[A]_S\le_S[B]_S\leftrightarrow A\preceq_S B}{\FormulaRefAuto{NullProfileComparison}[thm]{1,8,9}}
 \proofstepwidestar[1,2,3,8,9]{A\preceq_S B\leftrightarrow\Phi(A)\preceq_T\Phi(B)}{\FormulaRefAuto{NullProfileIsoPreorder}[thm]{1,2,3,8,9}}
 \proofstepwidestar[3]{\Phi:K_S\to K_T}{\FormulaRefAuto{SemigroupIsoFunction}[thm]{3}}
 \proofstepwidestar[3,8,9]{\Phi(A),\Phi(B)\in K_T}{\FormulaRefAuto{F\colon A\to B\dsep x\in A\vdash F(x)\in B}[thm]{12,8,9}}
 \proofstepwidestar[2,3,8,9]{[\Phi(A)]_T\le_T[\Phi(B)]_T\leftrightarrow\Phi(A)\preceq_T\Phi(B)}{\FormulaRefAuto{NullProfileComparison}[thm]{2,13}}
 \proofstepwidestar[1,2,3,8,9]{\psi([A]_S)=[\Phi(A)]_T\land\psi([B]_S)=[\Phi(B)]_T}{\FormulaRefAuto{NullProfileInducedMapValue}[thm]{1,2,3,8,9}}
 \proofstepwidestar[1,2,3,8,9]{[A]_S\le_S[B]_S\leftrightarrow\psi([A]_S)\le_T\psi([B]_S)}{\rIE{\FormulaRefAuto{a=b\vdash b=a}[thm]{\rAEa{15}},\FormulaRefAuto{a=b\vdash b=a}[thm]{\rAEb{15}},\FormulaRefAuto{P\leftrightarrow Q,R\leftrightarrow Q\vdash P\leftrightarrow R}[thm]{\FormulaRefAuto{P\leftrightarrow Q,Q\leftrightarrow R\vdash P\leftrightarrow R}[thm]{10,11},14}}}
 \proofstepwidestar[1,2,3,8,9]{p\le_Sq\leftrightarrow\psi(p)\le_T\psi(q)}{\rIE{\FormulaRefAuto{a=b\vdash b=a}[thm]{\rAEb{8}},\FormulaRefAuto{a=b\vdash b=a}[thm]{\rAEb{9}},16}}
 \proofstepwidestar[1,2,3,4]{p\le_Sq\leftrightarrow\psi(p)\le_T\psi(q)}{\rEE{6,[8]\vdots\rEE{7,[9]\vdots17}}}
\end{tabproofwide}

\Needspace{9\baselineskip}
\FormulaThmDeltaK[Die Profilabbildung ist bijektiv]{%
 H_S\dsep H_T\dsep I\vdash\psi:Q_S\bij Q_T
}{NullProfileIsoBijection}{\DeltaRow{Mengen}{S\dsep T\dsep p\dsep q\dsep A\dsep B}}
\Needspace{6\baselineskip}
\begin{tabproofsplitwide}
\Needspace{5\baselineskip}
\proofpartwide{Injektivität aus der reflektierten Ordnung}
 \proofstepwidestar[1]{H_S}{\rA}
 \proofstepwidestar[2]{H_T}{\rA}
 \proofstepwidestar[3]{I}{\rA}
 \proofstepwidestar[1,2,3]{\psi:Q_S\to Q_T}{\rAEa{\FormulaRefAuto{NullProfileInducedMapValue}[thm]{1,2,3}}}
 \proofstepwidestar[1]{\PartOrd{Q_S,\le_S}}{\FormulaRefAuto{NullProfilePartialOrder}[thm]{1}}
 \proofstepwidestar[2]{\PartOrd{Q_T,\le_T}}{\FormulaRefAuto{NullProfilePartialOrder}[thm]{2}}
 \proofstepwidestar[7]{p,q\in Q_S}{\rA}
 \proofstepwidestar[8]{\psi(p)=\psi(q)}{\rA}
 \proofstepwidestar[1,2,3,7]{\psi(p),\psi(q)\in Q_T}{\FormulaRefAuto{F\colon A\to B\dsep x\in A\vdash F(x)\in B}[thm]{4,7}}
 \proofstepwidestar[1,2,3,7]{\psi(p)\le_T\psi(p)}{\FormulaRefAuto{Partielle Ordnung}[def]{6,9}}
 \proofstepwidestar[1,2,3,7,8]{\psi(p)\le_T\psi(q)\land\psi(q)\le_T\psi(p)}{\rAI{\rIE{8,10},\rIE{8,10}}}
 \proofstepwidestar[1,2,3,7,8]{p\le_Sq\land q\le_Sp}{\FormulaRefAuto{NullProfileIsoOrder}[thm]{1,2,3,7,11}}
 \proofstepwidestar[1,2,3,7,8]{p=q}{\FormulaRefAuto{OrderAntisymmetry}[ax]{5,7,12}}
 \proofstepwidestar[1,2,3]{\psi:Q_S\inj Q_T}{\FormulaRefAuto{Injektive Funktion}[def]{4,\rUI{\rUI{\rRI{7,\rRI{8,13}}}}}}
\closeproofpartwide
\Needspace{5\baselineskip}
\proofpartwide{Surjektivität aus Repräsentanten}
 \setcounter{proofstepnr}{14}
 \proofstepwidestar[15]{q\in Q_T}{\rA}
 \proofstepwidestar[2]{\EqRel{K_T,\sim_T}}{\FormulaRefAuto{NullProfileEquivalence}[thm]{2}}
 \proofstepwidestar[2,15]{\exists B\in K_T\;q=[B]_T}{\FormulaRefAuto{QuotientRepresentative}[thm]{16,15}}
 \proofstepwidestar[18]{B\in K_T\land q=[B]_T}{\rA}
 \proofstepwidestar[3]{\Phi:K_S\bij K_T}{\FormulaRefAuto{SemigroupIsoBijectiveAxiom}[ax]{3}}
 \proofstepwidestar[3]{\Phi^{-1}:K_T\to K_S}{\FormulaRefAuto{InverseFunctionType}[thm]{19}}
 \proofstepwidestar[3,18]{\Phi^{-1}(B)\in K_S}{\FormulaRefAuto{F\colon A\to B\dsep x\in A\vdash F(x)\in B}[thm]{20,18}}
 \proofstepwidestar[1]{\EqRel{K_S,\sim_S}}{\FormulaRefAuto{NullProfileEquivalence}[thm]{1}}
 \proofstepwidestar[1,3,18]{[\Phi^{-1}(B)]_S\in Q_S}{\FormulaRefAuto{EqClassInQuotient}[thm]{22,21}}
 \proofstepwidestar[1,2,3,18]{\psi([\Phi^{-1}(B)]_S)=[\Phi(\Phi^{-1}(B))]_T}{\FormulaRefAuto{NullProfileInducedMapValue}[thm]{1,2,3,21}}
 \proofstepwidestar[3,18]{\Phi(\Phi^{-1}(B))=B}{\FormulaRefAuto{InverseFunctionRightInverse}[thm]{19,18}}
 \proofstepwidestar[1,2,3,18]{\psi([\Phi^{-1}(B)]_S)=q}{\rIE{25,\FormulaRefAuto{a=b\vdash b=a}[thm]{\rAEb{18}},24}}
 \proofstepwidestar[1,2,3,18]{\exists p\in Q_S\;\psi(p)=q}{\rEI{\rAI{23,26}}}
 \proofstepwidestar[1,2,3,15]{\exists p\in Q_S\;\psi(p)=q}{\rEE{17,[18]\vdots27}}
 \proofstepwidestar[1,2,3]{\psi:Q_S\sur Q_T}{\FormulaRefAuto{Surjektive Funktion}[def]{4,\rUI{\rRI{15,28}}}}
 \proofstepwidestar[1,2,3]{\psi:Q_S\bij Q_T}{\FormulaRefAuto{Bijektive Funktion}[def]{14,29}}
\closeproofpartwide
\end{tabproofsplitwide}

\Needspace{9\baselineskip}
\FormulaThmDeltaK[Bijektiver Transport der unteren Mengen]{%
 H_S\dsep H_T\dsep I\dsep p\in Q_S\vdash
 \Phi\restriction_{L_p}:L_p\bij L'_{\psi(p)}
}{NullProfileLowerTransport}{%
 \DeltaRow{Mengen}{S\dsep T\dsep p\dsep A\dsep L_p\dsep L'_{\psi(p)}}}
\begin{tabproofwide}
 \proofstepwidestar[1]{H_S}{\rA}
 \proofstepwidestar[2]{H_T}{\rA}
 \proofstepwidestar[3]{I}{\rA}
 \proofstepwidestar[4]{p\in Q_S}{\rA}
 \proofstepwidestar[3]{\Phi:K_S\bij K_T}{\FormulaRefAuto{SemigroupIsoBijectiveAxiom}[ax]{3}}
 \proofstepwidestar[1]{L_p\subseteq K_S}{\FormulaRefAuto{\{x\in A\mid P(x)\}\subseteq A}[thm]{\FormulaRefAuto{NullProfileSetsDef}[def]{1}}}
 \proofstepwidestar[2]{L'_{\psi(p)}\subseteq K_T}{\FormulaRefAuto{\{x\in A\mid P(x)\}\subseteq A}[thm]{\FormulaRefAuto{NullProfileSetsDef}[def]{2}}}
 \proofstepwidestar[8]{A\in K_S}{\rA}
 \proofstepwidestar[3,8]{\Phi(A)\in K_T}{\FormulaRefAuto{F\colon A\to B\dsep x\in A\vdash F(x)\in B}[thm]{\FormulaRefAuto{Bijektive Funktion}[def]{5},8}}
 \proofstepwidestar[1]{\EqRel{K_S,\sim_S}}{\FormulaRefAuto{NullProfileEquivalence}[thm]{1}}
 \proofstepwidestar[1,8]{[A]_S\in Q_S}{\FormulaRefAuto{EqClassInQuotient}[thm]{10,8}}
 \proofstepwidestar[1,2,3,8]{\psi([A]_S)=[\Phi(A)]_T}{\FormulaRefAuto{NullProfileInducedMapValue}[thm]{1,2,3,8}}
 \proofstepwidestar[1,2,3,4,8]{[A]_S\le_Sp\leftrightarrow\psi([A]_S)\le_T\psi(p)}{\FormulaRefAuto{NullProfileIsoOrder}[thm]{1,2,3,11,4}}
 \proofstepwidestar[1,8]{A\in L_p\leftrightarrow[A]_S\le_Sp}{\FormulaRefAuto{NullProfileSetsDef}[def]{1,8}}
 \proofstepwidestar[2,3,8]{\Phi(A)\in L'_{\psi(p)}\leftrightarrow[\Phi(A)]_T\le_T\psi(p)}{\FormulaRefAuto{NullProfileSetsDef}[def]{2,9}}
 \proofstepwidestar[1,2,3,4,8]{A\in L_p\leftrightarrow\Phi(A)\in L'_{\psi(p)}}{\FormulaRefAuto{P\leftrightarrow Q,R\leftrightarrow Q\vdash P\leftrightarrow R}[thm]{\FormulaRefAuto{P\leftrightarrow Q,Q\leftrightarrow R\vdash P\leftrightarrow R}[thm]{14,\rIE{12,13}},15}}
 \proofstepwidestar[1,2,3,4]{\forall A\in K_S\;(A\in L_p\leftrightarrow\Phi(A)\in L'_{\psi(p)})}{\rUI{\rRI{8,16}}}
 \proofstepwidestar[1,2,3,4]{\Phi\restriction_{L_p}:L_p\bij L'_{\psi(p)}}{\FormulaRefAuto{PointwiseCompatibilityImpliesBijectiveRestriction}[thm]{5,6,7,17}}
\end{tabproofwide}

\Needspace{9\baselineskip}
\FormulaThmDeltaK[Gleich große untere Grundelementmengen]{%
 H_S\dsep H_T\dsep I\dsep\Finite S\dsep\Finite T\dsep p\in Q_S\vdash
 \FiniteCard(D_p)=\FiniteCard(D'_{\psi(p)})
}{NullProfileLowerCardinality}{\DeltaRow{Mengen}{S\dsep T\dsep p\dsep D_p\dsep D'_{\psi(p)}}}
\begin{tabproofwide}
 \proofstepwidestar[1]{H_S}{\rA}
 \proofstepwidestar[2]{H_T}{\rA}
 \proofstepwidestar[3]{I}{\rA}
 \proofstepwidestar[4]{\Finite S}{\rA}
 \proofstepwidestar[5]{\Finite T}{\rA}
 \proofstepwidestar[6]{p\in Q_S}{\rA}
 \proofstepwidestar[1]{D_p\subseteq S}{\FormulaRefAuto{\{x\in A\mid P(x)\}\subseteq A}[thm]{\FormulaRefAuto{NullProfileSetsDef}[def]{1}}}
 \proofstepwidestar[2]{D'_{\psi(p)}\subseteq T}{\FormulaRefAuto{\{x\in A\mid P(x)\}\subseteq A}[thm]{\FormulaRefAuto{NullProfileSetsDef}[def]{2}}}
 \proofstepwidestar[1,4]{\Finite{D_p}}{\FormulaRefAuto{FiniteSubset}[thm]{7,4}}
 \proofstepwidestar[2,5]{\Finite{D'_{\psi(p)}}}{\FormulaRefAuto{FiniteSubset}[thm]{8,5}}
 \proofstepwidestar[1,2,3,6]{\Phi\restriction_{L_p}:L_p\bij L'_{\psi(p)}}{\FormulaRefAuto{NullProfileLowerTransport}[thm]{1,2,3,6}}
 \proofstepwidestar[1,2,3,6]{L_p\EqCard L'_{\psi(p)}}{\FormulaRefAuto{Gleichmächtigkeit}[def]{\rEI{11}}}
 \proofstepwidestar[1,6]{L_p=\PowNE{D_p}}{\FormulaRefAuto{NullProfileLowerPowerSet}[thm]{1,6}}
 \proofstepwidestar[1,2,3]{\psi:Q_S\to Q_T}{\rAEa{\FormulaRefAuto{NullProfileInducedMapValue}[thm]{1,2,3}}}
 \proofstepwidestar[1,2,3,6]{\psi(p)\in Q_T}{\FormulaRefAuto{F\colon A\to B\dsep x\in A\vdash F(x)\in B}[thm]{14,6}}
 \proofstepwidestar[1,2,3,6]{L'_{\psi(p)}=\PowNE{D'_{\psi(p)}}}{\FormulaRefAuto{NullProfileLowerPowerSet}[thm]{2,15}}
 \proofstepwidestar[1,2,3,6]{\PowNE{D_p}\EqCard\PowNE{D'_{\psi(p)}}}{\rIE{16,\rIE{13,12}}}
 \proofstepwidestar[1,2,3,4,5,6]{\FiniteCard(D_p)=\FiniteCard(D'_{\psi(p)})}{\FormulaRefAuto{FiniteNonemptyPowerSetReflectsCardinality}[thm]{9,10,17}}
\end{tabproofwide}

\Needspace{16\baselineskip}
\section{Fasergrößen und profiltreue Bijektionen}

\Needspace{12\baselineskip}
\FormulaDefDeltaK[Profilabbildung der Grundelemente]{%
 r_S:S\to Q_S,\qquad\forall s\in S\;(r_S(s)\coloneqq[\{s\}]_S)
}{NullProfileSingletonMapDef}{%
 \DeltaPrem{Halbgruppe mit Null}{H_S}
 \DeltaRow{Mengen}{S\dsep s\dsep Q_S}}
\begin{tabproof}
 \proofstep{1}{H_S}{\rA}
 \proofstep{2}{s\in S}{\rA}
 \proofstep{1}{\EqRel{K_S,\sim_S}}{\FormulaRefAuto{NullProfileEquivalence}[thm]{1}}
 \proofstep{2}{\{s\}\in K_S}{\FormulaRefAuto{SingletonPowerMember}[thm]{2}}
 \proofstep{1,2}{[\{s\}]_S\in Q_S}{\FormulaRefAuto{EqClassInQuotient}[thm]{3,4}}
\end{tabproof}

\Needspace{9\baselineskip}
\FormulaThmDeltaK[Die Profilmenge einer endlichen Halbgruppe ist endlich]{%
 H_S\dsep\Finite S\vdash\Finite{Q_S}
}{NullProfileQuotientFinite}{\DeltaRow{Mengen}{S\dsep Q_S\dsep p}}
\begin{tabproofwide}
 \proofstepwidestar[1]{H_S}{\rA}
 \proofstepwidestar[2]{\Finite S}{\rA}
 \proofstepwidestar[2]{\Finite{K_S}}{\FormulaRefAuto{FiniteNonemptyPowerSet}[thm]{2}}
 \proofstepwidestar[2]{\Finite{\powerset(K_S)}}{\FormulaRefAuto{FinitePowerSet}[thm]{3}}
 \proofstepwidestar[1]{\EqRel{K_S,\sim_S}}{\FormulaRefAuto{NullProfileEquivalence}[thm]{1}}
 \proofstepwidestar[6]{p\in Q_S}{\rA}
 \proofstepwidestar[1,6]{p\in\powerset(K_S)}{\rAEa{\FormulaRefAuto{QuotientSetChar}[thm]{5,6}}}
 \proofstepwidestar[1]{Q_S\subseteq\powerset(K_S)}{\FormulaRefAuto{SubsetIntroductionRule}[thm]{[6]\vdots7}}
 \proofstepwidestar[1,2]{\Finite{Q_S}}{\FormulaRefAuto{FiniteSubset}[thm]{8,4}}
\end{tabproofwide}

\Needspace{9\baselineskip}
\FormulaThmDeltaK[Profilfasern als Urbilder]{%
 H_S\dsep p\in Q_S\vdash C_p=r_S^{-1}[\{p\}]
}{NullProfileFiberSets}{\DeltaRow{Mengen}{S\dsep p\dsep s\dsep C_p}}
\begin{tabproofwide}
 \proofstepwidestar[1]{H_S}{\rA}
 \proofstepwidestar[2]{p\in Q_S}{\rA}
 \proofstepwidestar[1]{r_S:S\to Q_S}{\FormulaRefAuto{NullProfileSingletonMapDef}[def]{1}}
 \proofstepwidestar[2]{\{p\}\subseteq Q_S}{\FormulaRefAuto{a\in A\vdash\{a\}\subseteq A}[thm]{2}}
 \proofstepwidestar[1]{C_p\subseteq S}{\FormulaRefAuto{\{x\in A\mid P(x)\}\subseteq A}[thm]{\FormulaRefAuto{NullProfileSetsDef}[def]{1}}}
 \proofstepwidestar[6]{s\in S}{\rA}
 \proofstepwidestar[1,6]{r_S(s)=[\{s\}]_S}{\FormulaRefAuto{NullProfileSingletonMapDef}[def]{1,6}}
 \proofstepwidestar[1,6]{s\in C_p\leftrightarrow[\{s\}]_S=p}{\FormulaRefAuto{NullProfileSetsDef}[def]{1,6}}
 \proofstepwidestar[]{r_S(s)\in\{p\}\leftrightarrow r_S(s)=p}{\FormulaRefAuto{x\in\{a\}\eqvdash x=a}[thm]}
 \proofstepwidestar[1,6]{r_S(s)\in\{p\}\leftrightarrow s\in C_p}{\FormulaRefAuto{P\leftrightarrow Q,R\leftrightarrow Q\vdash P\leftrightarrow R}[thm]{9,\rIE{\FormulaRefAuto{a=b\vdash b=a}[thm]{7},8}}}
 \proofstepwidestar[1]{\forall s\in S\;(r_S(s)\in\{p\}\leftrightarrow s\in C_p)}{\rUI{\rRI{6,10}}}
 \proofstepwidestar[1,2]{r_S^{-1}[\{p\}]=C_p}{\FormulaRefAuto{PreimageEqFromPointwiseIff}[thm]{3,4,5,11}}
 \proofstepwidestar[1,2]{C_p=r_S^{-1}[\{p\}]}{\FormulaRefAuto{a=b\vdash b=a}[thm]{12}}
\end{tabproofwide}

Für \(p\in Q_S\) bezeichne \(I_p=\{q\in Q_S\mid q\le_Sp\}\)
die zugehörige untere Hauptmenge; \(I'_q\) wird entsprechend in \(Q_T\)
gebildet. Die Mengen \(D_p\) sind genau ihre Urbilder unter \(r_S\).

\Needspace{9\baselineskip}
\FormulaThmDeltaK[Untere Grundelementmengen als Urbilder]{%
 H_S\dsep p\in Q_S\vdash D_p=r_S^{-1}[I_p]
}{NullProfileLowerPreimage}{\DeltaRow{Mengen}{S\dsep p\dsep s\dsep I_p\dsep D_p}}
\begin{tabproofwide}
 \proofstepwidestar[1]{H_S}{\rA}
 \proofstepwidestar[2]{p\in Q_S}{\rA}
 \proofstepwidestar[1]{r_S:S\to Q_S}{\FormulaRefAuto{NullProfileSingletonMapDef}[def]{1}}
 \proofstepwidestar[]{I_p\subseteq Q_S}{\FormulaRefAuto{\{x\in A\mid P(x)\}\subseteq A}[thm]}
 \proofstepwidestar[1]{D_p\subseteq S}{\FormulaRefAuto{\{x\in A\mid P(x)\}\subseteq A}[thm]{\FormulaRefAuto{NullProfileSetsDef}[def]{1}}}
 \proofstepwidestar[6]{s\in S}{\rA}
 \proofstepwidestar[1,6]{r_S(s)\in Q_S}{\FormulaRefAuto{F\colon A\to B\dsep x\in A\vdash F(x)\in B}[thm]{3,6}}
 \proofstepwidestar[1,6]{r_S(s)\in I_p\leftrightarrow r_S(s)\le_Sp}{\FormulaRefAuto{x\in\{u\in A\mid P(u)\}\eqvdash x\in A\land P(x)}[thm]{7}}
 \proofstepwidestar[1,6]{s\in D_p\leftrightarrow r_S(s)\le_Sp}{\rIE{\FormulaRefAuto{a=b\vdash b=a}[thm]{\FormulaRefAuto{NullProfileSingletonMapDef}[def]{1,6}},\FormulaRefAuto{NullProfileSetsDef}[def]{1,6}}}
 \proofstepwidestar[1,6]{r_S(s)\in I_p\leftrightarrow s\in D_p}{\FormulaRefAuto{P\leftrightarrow Q,R\leftrightarrow Q\vdash P\leftrightarrow R}[thm]{8,9}}
 \proofstepwidestar[1]{\forall s\in S\;(r_S(s)\in I_p\leftrightarrow s\in D_p)}{\rUI{\rRI{6,10}}}
 \proofstepwidestar[1]{r_S^{-1}[I_p]=D_p}{\FormulaRefAuto{PreimageEqFromPointwiseIff}[thm]{3,4,5,11}}
 \proofstepwidestar[1]{D_p=r_S^{-1}[I_p]}{\FormulaRefAuto{a=b\vdash b=a}[thm]{12}}
\end{tabproofwide}

\Needspace{9\baselineskip}
\FormulaThmDeltaK[Korrespondierende Singleton-Profilfasern sind gleichmächtig]{%
 \begin{aligned}[t]
 &H_S\dsep H_T\dsep I\dsep\Finite S\dsep\Finite T\vdash{}\\[-2pt]
 &\forall p\in Q_S\,\exists u\;(u:C_p\bij C'_{\psi(p)})
 \end{aligned}
}{NullProfileFiniteFiberBijections}{\DeltaRow{Mengen}{S\dsep T\dsep p\dsep q\dsep C_p\dsep C'_{\psi(p)}}}
\Needspace{6\baselineskip}
\begin{tabproofsplitwide}
\Needspace{5\baselineskip}
\proofpartwide{Endliche geordnete Profilträger}
 \proofstepwidestar[1]{H_S}{\rA}
 \proofstepwidestar[2]{H_T}{\rA}
 \proofstepwidestar[3]{I}{\rA}
 \proofstepwidestar[4]{\Finite S}{\rA}
 \proofstepwidestar[5]{\Finite T}{\rA}
 \proofstepwidestar[1]{\PartOrd{Q_S,\le_S}}{\FormulaRefAuto{NullProfilePartialOrder}[thm]{1}}
 \proofstepwidestar[2]{\PartOrd{Q_T,\le_T}}{\FormulaRefAuto{NullProfilePartialOrder}[thm]{2}}
 \proofstepwidestar[1,4]{\Finite{Q_S}}{\FormulaRefAuto{NullProfileQuotientFinite}[thm]{1,4}}
 \proofstepwidestar[1]{r_S:S\to Q_S}{\FormulaRefAuto{NullProfileSingletonMapDef}[def]{1}}
 \proofstepwidestar[2]{r_T:T\to Q_T}{\FormulaRefAuto{NullProfileSingletonMapDef}[def]{2}}
 \proofstepwidestar[1,2,3]{\psi:Q_S\bij Q_T}{\FormulaRefAuto{NullProfileIsoBijection}[thm]{1,2,3}}
 \proofstepwidestar[12]{p,q\in Q_S}{\rA}
 \proofstepwidestar[1,2,3,12]{p\le_Sq\leftrightarrow\psi(p)\le_T\psi(q)}{\FormulaRefAuto{NullProfileIsoOrder}[thm]{1,2,3,12}}
 \proofstepwidestar[1,2,3]{\forall p,q\in Q_S\;(p\le_Sq\leftrightarrow\psi(p)\le_T\psi(q))}{\rUI{\rUI{\rRI{12,13}}}}
\closeproofpartwide
\Needspace{5\baselineskip}
\proofpartwide{Untere Urbilder und einzelne Fasern}
 \setcounter{proofstepnr}{14}
 \proofstepwidestar[15]{p\in Q_S}{\rA}
 \proofstepwidestar[1,2,3,15]{\psi(p)\in Q_T}{\FormulaRefAuto{F\colon A\bij B\dsep x\in A\vdash F(x)\in B}[thm]{11,15}}
 \proofstepwidestar[1,2,3,4,5,15]{\FiniteCard(D_p)=\FiniteCard(D'_{\psi(p)})}{\FormulaRefAuto{NullProfileLowerCardinality}[thm]{1,2,3,4,5,15}}
 \proofstepwidestar[1,15]{D_p=r_S^{-1}[I_p]}{\FormulaRefAuto{NullProfileLowerPreimage}[thm]{1,15}}
 \proofstepwidestar[1,2,3,15]{D'_{\psi(p)}=r_T^{-1}[I'_{\psi(p)}]}{\FormulaRefAuto{NullProfileLowerPreimage}[thm]{2,16}}
 \proofstepwidestar[1,2,3,4,5,15]{\FiniteCard(r_S^{-1}[I_p])=\FiniteCard(r_T^{-1}[I'_{\psi(p)}])}{\rIE{19,\rIE{18,17}}}
 \proofstepwidestar[1,2,3,4,5]{\forall p\in Q_S\;\FiniteCard(r_S^{-1}[I_p])=\FiniteCard(r_T^{-1}[I'_{\psi(p)}])}{\rUI{\rRI{15,20}}}
 \proofstepwidestar[1,2,3,4,5]{\forall p\in Q_S\,\exists u\;(u:r_S^{-1}[\{p\}]\bij r_T^{-1}[\{\psi(p)\}])}{\FormulaRefAuto{FinitePosetCumulativeTransportBijection}[thm]{6,7,8,4,5,9,10,11,14,21}}
 \proofstepwidestar[1,2,3,4,5,15]{\exists u\;(u:r_S^{-1}[\{p\}]\bij r_T^{-1}[\{\psi(p)\}])}{\rRE{\rUE{22},15}}
 \proofstepwidestar[1,15]{C_p=r_S^{-1}[\{p\}]}{\FormulaRefAuto{NullProfileFiberSets}[thm]{1,15}}
 \proofstepwidestar[1,2,3,15]{C'_{\psi(p)}=r_T^{-1}[\{\psi(p)\}]}{\FormulaRefAuto{NullProfileFiberSets}[thm]{2,16}}
 \proofstepwidestar[1,2,3,4,5,15]{\exists u\;(u:C_p\bij C'_{\psi(p)})}{\rIE{\FormulaRefAuto{a=b\vdash b=a}[thm]{24},\FormulaRefAuto{a=b\vdash b=a}[thm]{25},23}}
 \proofstepwidestar[1,2,3,4,5]{\forall p\in Q_S\,\exists u\;(u:C_p\bij C'_{\psi(p)})}{\rUI{\rRI{15,26}}}
\closeproofpartwide
\end{tabproofsplitwide}

\Needspace{9\baselineskip}
\FormulaThmDeltaK[Faserbijektionen in Urbildschreibweise]{%
 \begin{aligned}[t]
 &H_S\dsep H_T\dsep I\dsep\Finite S\dsep\Finite T\vdash{}\\[-2pt]
 &\forall p\in Q_S\,\exists u\;(u:r_S^{-1}[\{p\}]\bij r_T^{-1}[\{\psi(p)\}])
 \end{aligned}
}{NullProfileFinitePreimageFiberBijections}{\DeltaRow{Mengen}{S\dsep T\dsep p}}
\begin{tabproofwide}
 \proofstepwidestar[1]{H_S}{\rA}
 \proofstepwidestar[2]{H_T}{\rA}
 \proofstepwidestar[3]{I}{\rA}
 \proofstepwidestar[4]{\Finite S\land\Finite T}{\rA}
 \proofstepwidestar[1,2,3,4]{\forall p\in Q_S\,\exists u\;(u:C_p\bij C'_{\psi(p)})}{\FormulaRefAuto{NullProfileFiniteFiberBijections}[thm]{1,2,3,4}}
 \proofstepwidestar[6]{p\in Q_S}{\rA}
 \proofstepwidestar[1,2,3,4,6]{\exists u\;(u:C_p\bij C'_{\psi(p)})}{\rRE{\rUE{5},6}}
 \proofstepwidestar[1,6]{C_p=r_S^{-1}[\{p\}]}{\FormulaRefAuto{NullProfileFiberSets}[thm]{1,6}}
 \proofstepwidestar[1,2,3]{\psi:Q_S\to Q_T}{\rAEa{\FormulaRefAuto{NullProfileInducedMapValue}[thm]{1,2,3}}}
 \proofstepwidestar[1,2,3,6]{\psi(p)\in Q_T}{\FormulaRefAuto{F\colon A\to B\dsep x\in A\vdash F(x)\in B}[thm]{9,6}}
 \proofstepwidestar[1,2,3,6]{C'_{\psi(p)}=r_T^{-1}[\{\psi(p)\}]}{\FormulaRefAuto{NullProfileFiberSets}[thm]{2,10}}
 \proofstepwidestar[1,2,3,4,6]{\exists u\;(u:r_S^{-1}[\{p\}]\bij r_T^{-1}[\{\psi(p)\}])}{\rIE{11,\rIE{8,7}}}
 \proofstepwidestar[1,2,3,4]{\forall p\in Q_S\,\exists u\;(u:r_S^{-1}[\{p\}]\bij r_T^{-1}[\{\psi(p)\}])}{\rUI{\rRI{6,12}}}
\end{tabproofwide}

Die Fasern \(C_p\) können leer sein: Ein Profil kann von einer mehrgliedrigen
Teilmenge erzeugt werden, ohne das Profil irgendeiner Einermenge zu sein.
Die allgemeinen Fasersätze verlangen deshalb keine Surjektivität von
\(r_S\) oder \(r_T\). Sie verwenden das in Band~9 bereits eingeführte
Auswahlprinzip zum gemeinsamen Wählen der Faserbijektionen; für den
vorliegenden endlichen Indexbereich ist nur endliche Auswahl nötig.

\Needspace{9\baselineskip}
\FormulaThmDeltaK[Profiltreue Bijektionen erhalten die Nullproduktrelation]{%
 \begin{aligned}[t]
 &H_S\dsep H_T\dsep I\dsep f:S\bij T\dsep
 \forall s\in S\;(r_T(f(s))=\psi(r_S(s)))\dsep x,y\in S\vdash{}\\[-2pt]
 &xy=0_S\leftrightarrow f(x)f(y)=0_T
 \end{aligned}
}{NullProfileFaithfulBijectionZeroRelation}{\DeltaRow{Mengen}{S\dsep T\dsep x\dsep y}}
\begin{tabproofwide}
 \proofstepwidestar[1]{H_S}{\rA}
 \proofstepwidestar[2]{H_T}{\rA}
 \proofstepwidestar[3]{I}{\rA}
 \proofstepwidestar[4]{f:S\bij T}{\rA}
 \proofstepwidestar[5]{\forall s\in S\;(r_T(f(s))=\psi(r_S(s)))}{\rA}
 \proofstepwidestar[6]{x,y\in S}{\rA}
 \proofstepwidestar[4,6]{f(x),f(y)\in T}{\FormulaRefAuto{F\colon A\bij B\dsep x\in A\vdash F(x)\in B}[thm]{4,6}}
 \proofstepwidestar[6]{\{x\},\{y\}\in K_S}{\FormulaRefAuto{SingletonPowerMember}[thm]{6}}
 \proofstepwidestar[4,6]{\{f(x)\},\{f(y)\}\in K_T}{\FormulaRefAuto{SingletonPowerMember}[thm]{7}}
 \proofstepwidestar[3]{\Phi:K_S\to K_T}{\FormulaRefAuto{SemigroupIsoFunction}[thm]{3}}
 \proofstepwidestar[3,6]{\Phi(\{x\}),\Phi(\{y\})\in K_T}{\FormulaRefAuto{F\colon A\to B\dsep x\in A\vdash F(x)\in B}[thm]{10,8}}
 \proofstepwidestar[1,6]{r_S(x)=[\{x\}]_S\land r_S(y)=[\{y\}]_S}{\FormulaRefAuto{NullProfileSingletonMapDef}[def]{1,6}}
 \proofstepwidestar[2,4,6]{r_T(f(x))=[\{f(x)\}]_T\land r_T(f(y))=[\{f(y)\}]_T}{\FormulaRefAuto{NullProfileSingletonMapDef}[def]{2,7}}
 \proofstepwidestar[1,2,3,6]{\psi([\{x\}]_S)=[\Phi(\{x\})]_T\land\psi([\{y\}]_S)=[\Phi(\{y\})]_T}{\FormulaRefAuto{NullProfileInducedMapValue}[thm]{1,2,3,8}}
 \proofstepwidestar[5,6]{r_T(f(x))=\psi(r_S(x))\land r_T(f(y))=\psi(r_S(y))}{\rAI{\rRE{\rUE{5},\rAEa{6}},\rRE{\rUE{5},\rAEb{6}}}}
 \proofstepwidestar[1,2,3,4,5,6]{[\{f(x)\}]_T=[\Phi(\{x\})]_T\land[\{f(y)\}]_T=[\Phi(\{y\})]_T}{\rIE{\rAEa{12},\rAEb{12},\rAEa{13},\rAEb{13},\rAEa{14},\rAEb{14},15}}
 \proofstepwidestar[2]{\EqRel{K_T,\sim_T}}{\FormulaRefAuto{NullProfileEquivalence}[thm]{2}}
 \proofstepwidestar[1,2,3,4,5,6]{\{f(x)\}\sim_T\Phi(\{x\})\land\{f(y)\}\sim_T\Phi(\{y\})}{\FormulaRefAuto{EqClassRelFromEqual}[thm]{17,9,11,16}}
 \proofstepwidestar[1,2,3,4,5,6]{\{f(x)\}\{f(y)\}=z_T\leftrightarrow\Phi(\{x\})\Phi(\{y\})=z_T}{\FormulaRefAuto{NullProfileProductInvariant}[thm]{2,9,11,18}}
 \proofstepwidestar[1,2,3,6]{\{x\}\{y\}=z_S\leftrightarrow\Phi(\{x\})\Phi(\{y\})=z_T}{\FormulaRefAuto{NullProfileIsoZeroProduct}[thm]{1,2,3,8}}
 \proofstepwidestar[1,6]{\{x\}\{y\}=\{xy\}}{\FormulaRefAuto{SingletonPowerProduct}[thm]{\FormulaRefAuto{SemigroupWithZeroBaseAxiom}[ax]{1},6}}
 \proofstepwidestar[2,4,6]{\{f(x)\}\{f(y)\}=\{f(x)f(y)\}}{\FormulaRefAuto{SingletonPowerProduct}[thm]{\FormulaRefAuto{SemigroupWithZeroBaseAxiom}[ax]{2},7}}
 \proofstepwidestar[1,2,3,4,5,6]{\{xy\}=\{0_S\}\leftrightarrow\{f(x)f(y)\}=\{0_T\}}{\rIE{22,\rIE{21,\FormulaRefAuto{P\leftrightarrow Q,R\leftrightarrow Q\vdash P\leftrightarrow R}[thm]{20,19}}}}
 \proofstepwidestar[]{\{xy\}=\{0_S\}\leftrightarrow xy=0_S}{\FormulaRefAuto{\{a\}=\{b\}\eqvdash a=b}[thm]}
 \proofstepwidestar[]{\{f(x)f(y)\}=\{0_T\}\leftrightarrow f(x)f(y)=0_T}{\FormulaRefAuto{\{a\}=\{b\}\eqvdash a=b}[thm]}
 \proofstepwidestar[1,2,3,4,5,6]{xy=0_S\leftrightarrow f(x)f(y)=0_T}{\rLRS{25,\rLRS{24,23}}}
\end{tabproofwide}

In den folgenden Tabellen kürzen wir ausschließlich die ausgeschriebene
Eigenschaft
\[
 \mathcal K(f)\;\equiv\;
 f:S\bij T\ \land\ \forall s\in S\;(r_T(f(s))=\psi(r_S(s)))
\]
ab. Das ist die Faserverträglichkeit aus
\FormulaRefAuto{FiberwiseBijectionExists}[thm].

\Needspace{9\baselineskip}
\FormulaThmDeltaK[Rekonstruktion der endlichen Nullproduktrelation]{%
 \begin{aligned}[t]
 &H_S\dsep H_T\dsep I\dsep\Finite S\dsep\Finite T\vdash{}\\[-2pt]
 &\exists f\;\bigl(\mathcal K(f)\land
   \forall x,y\in S\;(xy=0_S\leftrightarrow f(x)f(y)=0_T)\bigr)
 \end{aligned}
}{NullProfileReconstruction}{\DeltaRow{Mengen}{S\dsep T\dsep x\dsep y\dsep p}}
\begin{tabproofwide}
 \proofstepwidestar[1]{H_S}{\rA}
 \proofstepwidestar[2]{H_T}{\rA}
 \proofstepwidestar[3]{I}{\rA}
 \proofstepwidestar[4]{\Finite S}{\rA}
 \proofstepwidestar[5]{\Finite T}{\rA}
 \proofstepwidestar[1]{r_S:S\to Q_S}{\FormulaRefAuto{NullProfileSingletonMapDef}[def]{1}}
 \proofstepwidestar[2]{r_T:T\to Q_T}{\FormulaRefAuto{NullProfileSingletonMapDef}[def]{2}}
 \proofstepwidestar[1,2,3]{\psi:Q_S\bij Q_T}{\FormulaRefAuto{NullProfileIsoBijection}[thm]{1,2,3}}
 \proofstepwidestar[1,2,3,4,5]{\forall p\in Q_S\,\exists u\;(u:r_S^{-1}[\{p\}]\bij r_T^{-1}[\{\psi(p)\}])}{\FormulaRefAuto{NullProfileFinitePreimageFiberBijections}[thm]{1,2,3,4,5}}
 \proofstepwidestar[1,2,3,4,5]{\exists f\;\mathcal K(f)}{\FormulaRefAuto{FiberwiseBijectionExists}[thm]{6,7,8,9}}
 \proofstepwidestar[11]{\mathcal K(f)}{\rA}
 \proofstepwidestar[12]{x,y\in S}{\rA}
 \proofstepwidestar[1,2,3,11,12]{xy=0_S\leftrightarrow f(x)f(y)=0_T}{\FormulaRefAuto{NullProfileFaithfulBijectionZeroRelation}[thm]{1,2,3,\rAEa{11},\rAEb{11},12}}
 \proofstepwidestar[1,2,3,11]{\forall x,y\in S\;(xy=0_S\leftrightarrow f(x)f(y)=0_T)}{\rUI{\rUI{\rRI{12,13}}}}
 \proofstepwidestar[1,2,3,11]{\exists f\;(\mathcal K(f)\land\forall x,y\in S\;(xy=0_S\leftrightarrow f(x)f(y)=0_T))}{\rEI{\rAI{11,14}}}
 \proofstepwidestar[1,2,3,4,5]{\exists f\;(\mathcal K(f)\land\forall x,y\in S\;(xy=0_S\leftrightarrow f(x)f(y)=0_T))}{\rEE{10,[11]\vdots15}}
\end{tabproofwide}

\Needspace{9\baselineskip}
\FormulaThmDeltaK[Zwei vorgeschriebene Werte bei der Nullproduktrekonstruktion]{%
 \begin{aligned}[t]
 &H_S\dsep H_T\dsep I\dsep\Finite S\dsep\Finite T\dsep
 a_0,a_1\in S\dsep b_0,b_1\in T\dsep{}\\[-2pt]
 &a_0\neq a_1\dsep b_0\neq b_1\dsep
 \Phi(\{a_0\})\sim_T\{b_0\}\dsep\Phi(\{a_1\})\sim_T\{b_1\}\vdash{}\\[-2pt]
 &\exists f\;\Bigl(\mathcal K(f)\land f(a_0)=b_0\land f(a_1)=b_1\\[-2pt]
 &\hspace{23mm}\land\forall x,y\in S\;(xy=0_S\leftrightarrow f(x)f(y)=0_T)\Bigr)
 \end{aligned}
}{NullProfileTwoPointReconstruction}{%
 \DeltaRow{Mengen}{S\dsep T\dsep a_0\dsep a_1\dsep b_0\dsep b_1\dsep x\dsep y\dsep p}}
\Needspace{6\baselineskip}
\begin{tabproofsplitwide}
\Needspace{5\baselineskip}
\proofpartwide{Verträglichkeit der vorgeschriebenen Werte}
 \proofstepwidestar[1]{H_S}{\rA}
 \proofstepwidestar[2]{H_T}{\rA}
 \proofstepwidestar[3]{I}{\rA}
 \proofstepwidestar[4]{\Finite S\land\Finite T}{\rA}
 \proofstepwidestar[5]{a_0,a_1\in S\land b_0,b_1\in T}{\rA}
 \proofstepwidestar[6]{a_0\neq a_1\land b_0\neq b_1}{\rA}
 \proofstepwidestar[7]{\Phi(\{a_0\})\sim_T\{b_0\}\land\Phi(\{a_1\})\sim_T\{b_1\}}{\rA}
 \proofstepwidestar[1]{r_S:S\to Q_S}{\FormulaRefAuto{NullProfileSingletonMapDef}[def]{1}}
 \proofstepwidestar[2]{r_T:T\to Q_T}{\FormulaRefAuto{NullProfileSingletonMapDef}[def]{2}}
 \proofstepwidestar[1,2,3]{\psi:Q_S\bij Q_T}{\FormulaRefAuto{NullProfileIsoBijection}[thm]{1,2,3}}
 \proofstepwidestar[1,2,3,4]{\forall p\in Q_S\,\exists u\;(u:r_S^{-1}[\{p\}]\bij r_T^{-1}[\{\psi(p)\}])}{\FormulaRefAuto{NullProfileFinitePreimageFiberBijections}[thm]{1,2,3,4}}
 \proofstepwidestar[5]{\{a_0\},\{a_1\}\in K_S\land\{b_0\},\{b_1\}\in K_T}{\FormulaRefAuto{SingletonPowerMember}[thm]{5}}
 \proofstepwidestar[3]{\Phi:K_S\to K_T}{\FormulaRefAuto{SemigroupIsoFunction}[thm]{3}}
 \proofstepwidestar[3,5]{\Phi(\{a_0\}),\Phi(\{a_1\})\in K_T}{\FormulaRefAuto{F\colon A\to B\dsep x\in A\vdash F(x)\in B}[thm]{13,12}}
 \proofstepwidestar[2]{\EqRel{K_T,\sim_T}}{\FormulaRefAuto{NullProfileEquivalence}[thm]{2}}
 \proofstepwidestar[2,3,5,7]{[\Phi(\{a_0\})]_T=[\{b_0\}]_T\land[\Phi(\{a_1\})]_T=[\{b_1\}]_T}{\FormulaRefAuto{EqClassEqualFromRel}[thm]{15,14,12,7}}
 \proofstepwidestar[1,2,3,5]{\psi([\{a_0\}]_S)=[\Phi(\{a_0\})]_T\land\psi([\{a_1\}]_S)=[\Phi(\{a_1\})]_T}{\FormulaRefAuto{NullProfileInducedMapValue}[thm]{1,2,3,12}}
 \proofstepwidestar[1,5]{r_S(a_0)=[\{a_0\}]_S\land r_S(a_1)=[\{a_1\}]_S}{\FormulaRefAuto{NullProfileSingletonMapDef}[def]{1,5}}
 \proofstepwidestar[2,5]{r_T(b_0)=[\{b_0\}]_T\land r_T(b_1)=[\{b_1\}]_T}{\FormulaRefAuto{NullProfileSingletonMapDef}[def]{2,5}}
 \proofstepwidestar[1,2,3,5,7]{r_T(b_0)=\psi(r_S(a_0))\land r_T(b_1)=\psi(r_S(a_1))}{\rAI{\rIE{\FormulaRefAuto{a=b\vdash b=a}[thm]{\rAEa{16}},\FormulaRefAuto{a=b\vdash b=a}[thm]{\rAEa{17}},\FormulaRefAuto{a=b\vdash b=a}[thm]{\rAEa{18}},\rAEa{19}},\rIE{\FormulaRefAuto{a=b\vdash b=a}[thm]{\rAEb{16}},\FormulaRefAuto{a=b\vdash b=a}[thm]{\rAEb{17}},\FormulaRefAuto{a=b\vdash b=a}[thm]{\rAEb{18}},\rAEb{19}}}}
\closeproofpartwide
\Needspace{5\baselineskip}
\proofpartwide{Faserweise Bijektion und Nullprodukte}
 \setcounter{proofstepnr}{20}
 \proofstepwidestar[1,2,3,4,5,6,7]{\exists f\;(\mathcal K(f)\land f(a_0)=b_0\land f(a_1)=b_1)}{\FormulaRefAuto{FiberwiseTwoPointBijectionExists}[thm]{8,9,10,11,5,6,20}}
 \proofstepwidestar[22]{\mathcal K(f)\land f(a_0)=b_0\land f(a_1)=b_1}{\rA}
 \proofstepwidestar[23]{x,y\in S}{\rA}
 \proofstepwidestar[1,2,3,22,23]{xy=0_S\leftrightarrow f(x)f(y)=0_T}{\FormulaRefAuto{NullProfileFaithfulBijectionZeroRelation}[thm]{1,2,3,\rAEa{22},23}}
 \proofstepwidestar[1,2,3,22]{\forall x,y\in S\;(xy=0_S\leftrightarrow f(x)f(y)=0_T)}{\rUI{\rUI{\rRI{23,24}}}}
 \proofstepwidestar[1,2,3,22]{\exists f\;\bigl(\mathcal K(f)\land f(a_0)=b_0\land f(a_1)=b_1\land\forall x,y\in S\;(xy=0_S\leftrightarrow f(x)f(y)=0_T)\bigr)}{\rEI{\rAI{22,25}}}
 \proofstepwidestar[1,2,3,4,5,6,7]{\exists f\;\bigl(\mathcal K(f)\land f(a_0)=b_0\land f(a_1)=b_1\land\forall x,y\in S\;(xy=0_S\leftrightarrow f(x)f(y)=0_T)\bigr)}{\rEE{21,[22]\vdots26}}
\closeproofpartwide
\end{tabproofsplitwide}

Die beiden markierten Elemente dürfen dasselbe Profil haben. In diesem
Fall werden beide Werte innerhalb derselben Faser festgelegt; die
Verschiedenheit der Ausgangs- und der Zielpunkte ist genau die benötigte
Voraussetzung des Zweipunktsatzes. Eine Erhaltung sämtlicher Einermengen
durch \(\Phi\) ist weder vorausgesetzt noch behauptet.
